% concatenated from Length_Maximal_nonlinear_codes_v3.tex
\documentclass[12pt, letterpaper]{article}

%%%% Standard Packages
\usepackage{amsmath,amssymb,amsfonts}
\usepackage{amsthm}
\usepackage{booktabs}
\usepackage{array}
\usepackage{float}
\usepackage[colorlinks=true,linkcolor=blue,citecolor=blue,urlcolor=blue]{hyperref}

%% Theorem Styles
\theoremstyle{plain}
\newtheorem{theorem}{Theorem}[section]
\newtheorem{proposition}[theorem]{Proposition}
\newtheorem{lemma}[theorem]{Lemma}
\newtheorem{corollary}[theorem]{Corollary}

\theoremstyle{definition}
\newtheorem{definition}[theorem]{Definition}
\newtheorem{example}[theorem]{Example}
\newtheorem{remark}[theorem]{Remark}

% Math Commands
\newcommand{\F}{\mathbb{F}}
\newcommand{\A}{{\mathcal{A}}}
\newcommand{\dist}{\operatorname{d}}
\newcommand{\mc}[1]{\mathcal{#1}}

\title{Length-Maximal Codes with Given Singleton Defect:\\ Structure and Bounds}

% FORMATTING  ----------------------------------------------------
\setlength{\textwidth}{6.5in}   % Change to 1-inch margins.
\setlength{\textheight}{10in}
\setlength{\evensidemargin}{0in} \setlength{\oddsidemargin}{0in}
\setlength{\topmargin}{-1.2in}
\setlength{\parindent}{0cm}
\setlength{\parsep}{1cm}

%\date{\today}

	\author{Tim Alderson\\Department of Mathematics and Statistics,\\ University of New Brunswick Saint John, \\ Saint John, NB, Canada. \thanks{The author acknowledges the support of the Natural Sciences and Engineering Research Council of Canada (NSERC), [funding reference number 2019-04103]\\
			Cette recherche a \'{e}t\'{e} financ\'{e}e par le Conseil de recherches en sciences naturelles et en g\'{e}nie du Canada (CRSNG), [num\'{e}ro de r\'{e}f\'{e}rence 2019-04103].} }

\begin{document}
	
	\maketitle

\begin{abstract}
	Let $\mathcal C$ be an $(n,q^k,d)_q$ code of integer dimension $k\ge2$ and Singleton defect $s=n-k+1-d$. Repeated shortening followed by the Plotkin bound gives
	\[
	 n\le (s+1)(q+1)+k-2.
	\]
	We study the equality case, calling a code attaining this bound \emph{length-maximal}. Equality provides rigid structure: the code is symbol-uniform under every successive shortening, is an orthogonal array of strength at least $k-1$, and has pairwise distances in $\{d\}\cup\{n-k+3,\ldots,n\}$. For $k\ge3$ it also satisfies $s\le q-1$ and $(s+2)\mid q(q+1)$. In dimension two, length-maximal codes are equivalent to resolvable $2$-$(q^2,q,s+1)$ multidesigns.

	The (inner) distance distribution, and hence the Hamming weight enumerator after any codeword is normalised to zero, is determined by the parameters. In dimension four this gives three cases: apart from the MDS case, only $s=q-2$ and $s=q-1$ remain possible; in the latter case the alphabet size is constrained by $(q+2)\mid36$, leaving six values of $q$. For $q\ge4$, every length-maximal code of dimension at least five is MDS, and requires $q$ to be a multiple of $36$. The non-MDS cases are the dual and extended ternary Golay codes, of dimensions five and six. Consequently, for $s\ge1$ one has $k\le4$ unless $q=3$, where $k\le6$.

	By shortening the codes before applying Plotkin, we provide sharper large-defect bounds and several ranges in which nonlinear codes satisfy the Griesmer bound.
\end{abstract}

\medskip

	\noindent\textbf{Keywords:} nonlinear codes; Singleton defect; Plotkin bound; orthogonal arrays; maximal arcs; resolvable designs.

\smallskip

\noindent\textbf{MSC 2020:} 94B65; 94B25; 94B60; 05B05.

	\newpage
	%==============================================================================
	% SECTION 1: INTRODUCTION
	%==============================================================================
	\section{Introduction}
	
	An \emph{$(n,M,d)_q$ code} $\mc C$ is an $M$-subset of $\A^n$, where $|\A|=q$ and $d$ is the minimum Hamming distance between distinct codewords. If $\A=\F_q$ and $\mc C$ is a $k$-dimensional subspace of $\F_q^n$, it is a linear $[n,k,d]_q$ code. More generally, the dimension of an arbitrary code is $\log_q M$; throughout the main results we assume this dimension is an integer, so $M=q^k$. The \emph{Singleton defect} of
	$\mathcal{C}$ is
	\[
	s = n - k + 1 - d \ge 0,
	\]
	measuring how far the code falls short of the Singleton bound \cite{Singleton1964, MacWilliamsSloane1977}. A code with $s = 0$ is \emph{maximum distance separable} (MDS); a code with $s = 1$ is \emph{almost maximum distance separable} (AMDS) \cite{DeBoer1996}; and more generally a code with defect $s$ is called an $A^s$MDS code.
	
	For linear codes, length questions are naturally geometric: the columns of a generator matrix form a multiset of points, a \emph{projective system}, in
	$\mathrm{PG}(k-1,q)$. For a fixed defect $s$, the length $n$ is maximised when this multiset is a maximal $(n, k+s-1)$-arc, that is, an arc meeting the upper bound
	\begin{equation}\label{eq:arc_bound}
		n \le (s+1)(q+1) + k - 2.
	\end{equation}
	Maximal arcs are rare. In the planar case ($k=3$), the celebrated theorem of Ball, Blokhuis, and Mazzocca \cite{Ball1997MaximalArcs} establishes that non-trivial maximal arcs (those with $s < q-2$) do not exist in $\mathrm{PG}(2,q)$ for odd $q$. In even characteristic, maximal arcs exist whenever $(s+2) \mid q$, due to constructions of Denniston \cite{Denniston1969}. 
	
	Although maximal arcs do not exist in Desarguesian planes of odd order, De Clerck, Delanote, Hamilton, and Mathon constructed a \emph{perp-system} of $21$ lines in $\mathrm{PG}(5,3)$ such that every hyperplane contains $0$ or $3$ of the lines \cite{DeClerck2002PerpSystems}. It gives an additive, non-$\F_9$-linear, length-maximal $(21,9^3,18)_9$ AMDS code \cite{alderson2026setssubspacesrestrictedhyperplane}. To the best of our knowledge, this is the only known non-trivial three-dimensional length-maximal code over an odd-order alphabet for parameters forbidden in the linear setting.

	A second question is which classical linear bounds persist without linearity. The Griesmer bound does not hold for arbitrary nonlinear codes. Guerrini, Meneghetti, and Sala \cite{Guerrini2015Optimal}, building on Bellini and Meneghetti \cite{Bellini2015Griesmer}, studied \emph{systematic} codes, that is, images of injections of the form $x\mapsto(x,f_{k+1}(x),\ldots,f_n(x))$. Adopting their notation, let $N_q(M,d)$, $S_q(k,d)$, and $L_q(k,d)$ be the minimum lengths of unrestricted, systematic, and linear codes, respectively. They proved, among other results, that the Griesmer bound holds for unrestricted codes when $M\ge q^k$ and $q^{k-1}\mid d$, and established for binary systematic codes the inequality
		\begin{equation}\label{eq:GMS_bound}
		S_2(k,d) \ge k + \left\lceil \tfrac{3}{2}d \right\rceil - 2.
	\end{equation}
	The Plotkin bound applies without linearity. Followed by repeated shortening, it gives \eqref{eq:arc_bound} for every $(n,q^k,d)_q$ code. The inequality itself is quite elementary, and the subject of this paper is the rigidity imposed by equality. We call a code attaining \eqref{eq:arc_bound} \emph{length-maximal}. The terminology reflects the connection with maximal arcs while allowing an arbitrary alphabet and no additivity assumption.
	
	Equality forces symbol-uniformity after every successive shortening. Equivalently, a length-maximal code is an orthogonal array of strength at least $k-1$. Its pairwise distances lie in $\{d\}\cup\{n-k+3,\ldots,n\}$, and for $k\ge3$ one has $s\le q-1$ and $(s+2)\mid q(q+1)$. Dimension-two codes correspond to resolvable $2$-$(q^2,q,s+1)$ multidesigns, a defect-based formulation of the classical correspondence of Semakov and Zinov'ev \cite{SemakovZinoviev1968} between equidistant codes meeting the Plotkin bound and resolvable designs. In dimension three the zero-sets of minimum-weight words form a $2$-design; a support of multiplicity $q-1$ forces the stronger linear divisibility $(s+2)\mid q$ when $s\le q-2$.

	The orthogonal-array moments determine the inner distance distribution from every codeword. For $k=4$ they yield three cases: no length-maximal code exists for $1\le s\le q-3$; when $s=q-2$ there is no full-distance pair; and when $s=q-1$ such pairs are forced and integrality requires $(q+2)\mid36$. For $q\ge4$, every length-maximal code of dimension at least five is MDS, giving nonexistence when $36\nmid q$. The exceptional non-MDS cases are the dual and extended ternary Golay codes. Thus the ordinary Hamming weight enumerator, after any codeword is normalised to zero, is determined by the parameters, although existence remains open in several families.
	
	Table~\ref{tab:by_dimension} summarises the existence picture. Outside the unresolved MDS case $36\mid q$, the known separation between the linear and unrestricted settings occurs only in dimensions two through four. In dimensions at least five and with $36\nmid q$, every length-maximal code is equivalent to a linear code.

	\begin{table}[htbp]
		\centering
		\small
		\renewcommand{\arraystretch}{1.25}
		\caption{Length-maximal codes by dimension $k$: where linear and genuinely nonlinear examples are known to exist, representative codes, and the relevant results. Rows account for all defects $s \ge 0$; for $q \ge 4$ and $k \ge 5$ every length-maximal code is necessarily MDS (Theorem~\ref{thm:dim5}), and for $q > 2$ and $k \ge 4$ the MDS case $s=0$ requires $36 \mid q$ (Theorem~\ref{thm:bruen_silverman}(2)), with no example known. For $k \ge 5$ and $36 \nmid q$, every length-maximal code is equivalent to a linear one (Corollary~\ref{cor:dim5_linearity}). The forced  distance distributions are given in Table~\ref{tab:enumerators}.}
		\label{tab:by_dimension}
		\begin{tabular}{@{}c >{\raggedright\arraybackslash}p{1.4in} >{\raggedright\arraybackslash}p{1.7in} >{\raggedright\arraybackslash}p{1.35in} >{\raggedright\arraybackslash}p{1in}@{}}
			\toprule
			$k$ & linear length-maximal & genuinely nonlinear & representative example(s) & references \\
			\midrule
			$2$ & all $s$ ($q$ a prime power) &  all $s$, whenever a non-Desarguesian affine plane of order $q$ exists (least $q{=}9$) & non-Desarguesian planes & Prop.~\ref{prop:2d_design_equiv}, Cor.~\ref{cor:2d_existence} \\
			$3$ & $s{\in}\{q{-}2,\,q{-}1\}$ (all $q$); even $q$ with $(s{+}2)\mid q$; none for odd $q$, $0\le s\le q{-}3$ & $s{=}1$, $q{=}9$; possible only if $(s{+}2)\mid q(q{+}1)$ & perp-system $(21,9^3,18)_9$ & Lem.~\ref{lem:divisibility}, Thm.~\ref{thm:zero_set_design}, Cor.~\ref{cor:dim23_enumerator}; \cite{DeClerck2002PerpSystems} \\
			$4$ & $s{=}q{-}2$ (all prime powers $q$); $s{=}q{-}1$ only $q{=}2$ & none known. Possible only if $q>2$ and: $s{=}0$ with $36\mid q$; $s{=}q{-}1$ with $(q{+}2)\mid36$; or $s{=}q{-}2$ & ovoid codes; extended Hamming $[8,4,4]_2$ & Thm.~\ref{thm:dim4_enumerator}, Cor.~\ref{cor:dim4_trichotomy}, Rem.~\ref{rem:dim4_q1_candidates} \\
			$5$ & $q{=}2$, $s{=}0$; $q{=}3$, $s{=}1$ (both unique) & none known. Possible only if $36\mid q$ & dual Golay $[11,5,6]_3$ & Lem.~\ref{lem:dim5_enumerator}, Thm.~\ref{thm:dim5existence}, Cor.~\ref{cor:dim5_linearity} \\
			$6$ & $q{=}2$, $s{=}0$; $q{=}3$, $s{=}1$ (both unique) & none known. Possible only if $36\mid q$ & extended Golay $[12,6,6]_3$ & Prop.~\ref{prop:dim6_q3}, Rem.~\ref{rem:uniqueness}, Cor.~\ref{cor:dim5_linearity} \\
			$\ge 7$ & only $q{=}2$, $s{=}0$ (unique) & none known. Possible only if $36\mid q$ & even-weight $(k{+}1,2^k,2)_2$ & Prop.~\ref{prop:q3_dimension}, Cor.~\ref{cor:dim5_linearity} \\
			\bottomrule
		\end{tabular}
	\end{table}

	For general, not necessarily length-maximal codes, shortening to any intermediate dimension and applying Plotkin gives a hierarchy of length bounds (Theorem~\ref{thm:plotkin_shortening_hierarchy}). Its first member is \eqref{eq:arc_bound}, with subsequent members offering improvement as the defect grows. The hierarchy and the binary Plotkin refinements also identify several ranges in which unrestricted nonlinear codes satisfy the Griesmer bound, collected in Table~\ref{tab:griesmer_conditions}.
	
Section~\ref{sec:prelim} fixes notation and recalls the required bounds. Section~\ref{sec:maximal_arc_bound} develops the equality theory and the dimension-by-dimension classification. Section~\ref{sec:refinements} gives binary refinements, the general Plotkin--shortening hierarchy, and Griesmer consequences. Section~\ref{sec:conclusion} concludes with open problems.

\subsection{Notation}

We collect the main notation and standing assumptions used throughout the paper in Table~\ref{tab:notation}.

\begin{table}[htbp]
	\centering
	\renewcommand{\arraystretch}{1.2}
	\begin{tabular}{@{}l p{4.15in}@{}}
		\toprule
		\textbf{Notation} & \textbf{Meaning} \\
		\midrule
		$(n,M,d)_q$ code $\mathcal{C}$ 
		& $M$-element subset of $\A^n$ ($|\A|=q$) with minimum Hamming distance exactly $d$ \\
		$\dist(\mathbf{u},\mathbf{v})$ 
		& Hamming distance between codewords $\mathbf{u}$ and $\mathbf{v}$ \\
		$k = \log_q M$
		& dimension of the code; throughout, $M = q^k$ \\
		$s = n - k + 1 - d$
		& Singleton defect ($\ge 0$); the code is called an $A^s$MDS code \\
		$\mathbf{0} \in \mathcal{C}$ 
		& standing assumption (without loss of generality, by symbol permutation) \\
		symbol-uniform 
		& every coordinate has each symbol of $\A$ appearing exactly $M/q$ times (requires $q \mid M$) \\
		$\mathcal{C}_i^{(a)}$ 
		& shortened code obtained by deleting coordinate $i$ from all codewords with $i$-th entry equal to $a$ \\
		$g_q(k,d) = \sum_{i=0}^{k-1} \lceil d/q^i \rceil$ 
		& Griesmer bound for linear $[n,k,d]_q$ codes \\
		\midrule
		\multicolumn{2}{@{}p{6.3in}@{}}{\textit{Standing assumptions}} \\
		\multicolumn{2}{@{}p{6.3in}@{}}{\quad $k \ge 2$ for all main results in Sections \ref{sec:maximal_arc_bound} and~\ref{sec:refinements}} \\
		\multicolumn{2}{@{}p{6.3in}@{}}{\quad A code is \emph{length-maximal} if it attains the bound of Theorem \ref{thm:main_bound} } \\
		\bottomrule
	\end{tabular}
	\caption{Notation and standing assumptions.}
	\label{tab:notation}
\end{table}
	%==============================================================================
	% SECTION 2: PRELIMINARIES
	%==============================================================================
	\section{Preliminaries}\label{sec:prelim}
	
	In this section we establish notation, recall the fundamental bounds used 	throughout, and introduce the two key structural concepts of the paper: 	symbol-uniformity and shortening.
	
	\subsection{Code Equivalence}

	Two codes $\mathcal{C}, \mathcal{C}' \subseteq \A^n$ are \emph{equivalent} if one can be obtained from the other by a combination of \emph{positional permutations} (permuting the $n$ coordinate positions) and \emph{symbol permutations} (applying, independently in each coordinate, a bijection of $\A$) \cite{Welsh1988}. Equivalent codes have identical parameters $n$, $M$, $d$. Since an appropriate symbol permutation can map any chosen codeword to the all-zero word, we may assume without loss of generality that $\mathbf{0} \in \mathcal{C}$, and we adopt this convention throughout. A code $\mc{C}$ is said to be \textit{equivalent to linear} if there is a linear code $\mc{C}'$ that is equivalent to $\mc{C}$; a code that is not equivalent to linear is called \emph{genuinely nonlinear}.
	
	\subsection{Fundamental Bounds}
	
	The Singleton bound holds for all codes, regardless of linearity or structure.
	
	\begin{theorem}[Singleton Bound \cite{Singleton1964, MacWilliamsSloane1977}]
		For any $(n, M, d)_q$ code with dimension $\kappa = \log_q M$,
		\[
		d \le n - \lceil \kappa \rceil + 1.
		\]
	\end{theorem}
	
	\begin{definition}
		The \emph{Singleton defect} of an $(n, M, d)_q$ code $\mathcal{C}$ is
		\[
		s = n - \lceil \kappa \rceil + 1 - d \ge 0.
		\]
		We recall that a code with defect $  s  $ is an $A^s$MDS code. When $  M=q^k  $, this reduces to the standard linear defect $  s=n-k+1-d  $.
	\end{definition}

	The Plotkin bound gives an upper bound on the size of a code when the minimum distance is large relative to the length. The following $q$-ary form underlies both the main length bound and the hierarchy of Section~\ref{sec:large_defect}.
	
	\begin{theorem}[$q$-ary Plotkin Bound \cite{Plotkin1960,SilvermanMetrization1960,MR238597}]
		\label{thm:plotkin_qary}
		Every $(n,M,d)_q$ code satisfies
		\begin{equation}\label{eq:plotkin_length}
		 n\ge \left\lceil d\,\frac{1-1/M}{1-1/q}\right\rceil.
		\end{equation}
		If, in addition, $dq>n(q-1)$, then
		\[
		 M\le \left\lfloor\frac{dq}{dq-n(q-1)}\right\rfloor.
		\]
	\end{theorem}
	
	For binary codes, a sharper formulation that distinguishes the parity of $d$ will be needed in Section~\ref{sec:binary_refinements}.
	
	\begin{theorem}[Binary Plotkin Bound
		\cite{Plotkin1960, MacWilliamsSloane1977}]
		\label{thm:plotkin_binary}
		Let $\mathcal{C}$ be a binary $(n, M, d)_2$ code.
		\begin{enumerate}
			\item If $d$ is even and $2d > n$, then $M \le 2\left\lfloor \dfrac{d}{2d-n} \right\rfloor$.
			\item If $d$ is odd and $2d+1 > n$, then $M \le 2\left\lfloor \dfrac{d+1}{2d+1-n} \right\rfloor$.
		\end{enumerate}
	\end{theorem}
	
	For linear codes, the Griesmer bound provides a stronger lower bound on length than the Singleton bound.
	
	\begin{theorem}[Griesmer Bound \cite{Griesmer1960, SS65}]
		\label{thm:griesmer}
		For any linear $[n, k, d]_q$ code,

		\begin{equation} \label{eqn:GriesmerBound}
			n \ge g_q(k,d) := \sum_{i=0}^{k-1} \left\lceil \frac{d}{q^i} \right\rceil.
		\end{equation}

	\end{theorem}

	The proof of the Griesmer bound relies critically on linearity, and it is known that the bound does not hold for general nonlinear codes \cite{Bellini2015Griesmer}. One aim of this paper is to identify conditions under which nonlinear codes of fixed Singleton defect respect this bound.

	Finally, we record a nonexistence theorem of Bruen and Silverman for MDS codes of maximal length, which will be invoked repeatedly in
	Section~\ref{sec:maximal_arc_bound}. In the notation of \cite{BruenSilverman1983}, an $L(n,k,r)$ space is a code of size $n^r$ and 	length $k$ over an \emph{arbitrary} alphabet of $n$ symbols in which no two codewords agree in $r$ or more positions; in our notation this is precisely an $(n', q^{k'}, d)_q$ MDS code with $q = n$, $n' = k$, and $k' = r$.  Every MDS code satisfies $n' \le q + k' - 1$ 	\cite[Lemma~5]{SilvermanMetrization1960}, the case $s=0$ of Theorem~\ref{thm:main_bound}. The following theorem concerns the equality case, that is, length-maximal MDS codes.

	\begin{theorem}[Bruen--Silverman {\cite[Theorems~4 and~6]{BruenSilverman1983}}]
		\label{thm:bruen_silverman}
		Let $q > 2$ and $k \ge 3$ be integers, and suppose a $(q+k-1,\, q^k,\, q)_q$ MDS code exists.
		\begin{enumerate}
			\item Then $4 \mid q$; in particular no such code exists for odd $q$.
			\item If moreover $k \ge 4$, then $36 \mid q$.
		\end{enumerate}
	\end{theorem}

	\begin{remark}\label{rem:bs_nonlinear}
		Theorem~\ref{thm:bruen_silverman} holds for every integer $q > 2$, and the hypothetical extremal codes of part~(2) are necessarily genuinely 	nonlinear (no prime power is divisible by $36$). Whether any code attaining the hypotheses of part~(2) exists remains open. The smallest candidate parameters are $(39,\,36^4,\,36)_{36}$.
	\end{remark}

	\subsection{Symbol-Uniformity}
	
	A key structural property of optimal codes, which will emerge naturally from our main bound, is the following.
	
	\begin{definition}
		A code $\mathcal{C} \subseteq \A^n$ is \emph{symbol-uniform in 	coordinate $j$} if every symbol $a \in \A$ appears equally often, that is,
		\[
		|\{ \mathbf{c} \in \mathcal{C} \mid c_j = a \}| = \frac{M}{q}
		\quad \text{for each } a \in \A.
		\]
		If $\mathcal{C}$ is symbol-uniform in every coordinate, it is called \emph{symbol-uniform}.
	\end{definition}
	
	Note that symbol-uniformity in any coordinate requires $q \mid M$.

	\subsection{Orthogonal Arrays}

	An $N\times n$ array over a $q$-symbol alphabet is an $\operatorname{OA}(N,n,q,t)$ of strength $t$ if, in every choice of $t$ columns, each ordered $t$-tuple occurs exactly $N/q^t$ times \cite{HedayatSloaneStufken1999}. Thus symbol-uniformity is precisely strength one. We shall show that a length-maximal code of dimension $k$ is an orthogonal array of strength at least $k-1$ (Corollary~\ref{cor:orthogonal_array}). This observation
	is leveraged later to determine the inner distance distribution.

	\subsection{Shortening}
	
	A standard operation that preserves many code properties is shortening. We 	use it extensively in our inductive arguments.
	
	\begin{definition}
		Let $\mathcal{C}$ be an $(n, M, d)_q$ code. For a coordinate $i \in \{1, \dots, n\}$ and a symbol $a \in \A$, the \emph{shortened code} is
		\[
		\mathcal{C}_i^{(a)} = \bigl\{ (c_1, \dots, c_{i-1}, c_{i+1}, \dots, c_n) \mid \mathbf{c} \in \mathcal{C},\ c_i = a \bigr\}.
		\]
	\end{definition}
	\begin{lemma}[Properties of Shortening]\label{lem:shortening}
		Let $\mathcal{C}$ be an $(n, M, d)_q$ $A^s$MDS code with $q^k \le M < q^{k+1}$, and let $\mathcal{C}_i^{(a)}$ be an $(n-1, M', d')_q$ $A^{s'}$MDS code of dimension $\kappa'$.
		\begin{enumerate}
			\item $M' = |\{ \mathbf{c} \in \mathcal{C} \mid c_i = a \}|$, and $d' \ge d$.
			\item The symbol $a$ can be chosen so that 	$M' \ge \lceil M/q \rceil$; consequently there is an $(n-1,\lceil M/q \rceil,d')$ A$^{s'}$MDS code $\mc{C}'\subseteq \mc{C}_i^{(a)}$, with $d'\ge d$ and  $s' \le s$.
			\label{item:shortening_size}
			\item If $\mathcal{C}$ is symbol-uniform in coordinate $i$, then $M' = M/q$.
			\item If $\mathcal{C}$ is symbol-uniform, then $a$ and $i$ can be chosen so that $s' = s$ (and hence $d' = d$).
		\end{enumerate}
	\end{lemma}

	\begin{proof}
		Items (1)--(3) follow directly from the definitions and simple counting arguments. For item (4), let $\mathbf{x}, \mathbf{y} \in \mathcal{C}$ be a pair at minimum distance $d$, and let $i$ be a coordinate in which $x_i = y_i$; such a coordinate exists, since $d \le n - k + 1 < n$ when $k \ge 2$. Shortening at coordinate $i$ with $a = x_i$ retains both $\mathbf{x}$ and $\mathbf{y}$, whose distance is unchanged, so $d' = d$ and hence $s' = s$.
	\end{proof}
	
	With these tools in place, we turn to the main bound of the paper.
	
	%==============================================================================
	% SECTION 3: THE MAXIMAL-ARC-TYPE BOUND
	%==============================================================================
	\section{Length-Maximal Codes}
	\label{sec:maximal_arc_bound}
	
	In this section we prove our main upper bound on the length of a code in terms of its alphabet size $q$, its (integer) dimension $k = \log_q M$, and its Singleton defect $s$. We then show that codes achieving this bound possess a rigid combinatorial structure: they are necessarily symbol-uniform and 	their pairwise distances are sharply constrained.
	
	\subsection{The Base Case: Codes of Dimension 2}
	\label{sec:base_case}
	
	We begin with codes of size $M = q^2$.
	
	\begin{lemma}\label{lem:2d_bound}
		If $\mathcal{C}$ is an $(n, q^2, d)_q$ code with Singleton defect 	$s = n - d - 1$, then
		\[
		n \le (s+1)(q+1).
		\]
		Moreover, equality holds if and only if $\mathcal{C}$ is symbol-uniform and equidistant, that is, every pair of distinct codewords agrees in
		exactly $s+1$ coordinates.
	\end{lemma}
	
	\begin{proof}
		Consider the complete graph on the $q^2$ codewords of $\mathcal{C}$. For distinct codewords $\mathbf{x}, \mathbf{y}$, let $a(\mathbf{x},\mathbf{y})$ denote the number of coordinates in which they agree. Since any two distinct codewords agree in at most $n-d = s+1$ positions,
		\[
		\sum_{\{\mathbf{x},\mathbf{y}\}} a(\mathbf{x},\mathbf{y}) \le \binom{q^2}{2}(s+1).
		\]
		
		On the other hand, count the agreements coordinate-wise. For each coordinate $i$, let $n_{i,a}$ be the number of codewords with symbol $a\in\A$ in position $i$, so $\sum_a n_{i,a}=q^2$. The contribution of coordinate $i$ is
		\[
		A_i = \sum_{a\in\A} \binom{n_{i,a}}{2}.
		\]
		By convexity of $\binom{\cdot}{2}$, $A_i$ is minimised when $n_{i,a}=q$ for all $a$, yielding
		\[
		A_i \ge q\binom{q}{2} = \frac{q^2(q-1)}{2},
		\]
		with equality if and only if $\mathcal{C}$ is symbol-uniform in coordinate $i$.
		
		Summing over all $n$ coordinates and combining, we obtain
		\begin{equation}\label{eqn:2d_double_count}
			n \cdot \frac{q^2(q-1)}{2} \le \sum_{i=1}^n A_i \le \frac{q^2(q^2-1)}{2}(s+1).
		\end{equation}
		
		This immediately yields $n \le (s+1)(q+1)$.
		
		Equality in \eqref{eqn:2d_double_count} forces $A_i = \frac{q^2(q-1)}{2}$ for every $i$, which by convexity requires symbol-uniformity in every coordinate. It also forces $a(\mathbf{x},\mathbf{y}) = s+1$ for every pair of distinct codewords (equidistance). Conversely, if $\mathcal{C}$ is symbol-uniform and equidistant with agreement count exactly $s+1$, then equality holds throughout \eqref{eqn:2d_double_count}, and thus $n = (s+1)(q+1)$.
	\end{proof}
	
	\begin{corollary}\label{cor:2d_LM_GB}
	Let $q$ be a prime power. With $N_q(q^2,d)$, $S_q(2,d)$, and $L_q(2,d)$ as defined in the Introduction, $N_q(q^2,d)=S_q(2,d)=L_q(2,d)= d+\left\lceil \frac{d}{q}\right\rceil=g_q(2,d)$.
	\end{corollary}
	\begin{proof}
		From Lemma~\ref{lem:2d_bound}, every $(n,q^2,d)_q$ $A^s$MDS code satisfies $n\le (s+1)(q+1)=(n-d)(q+1)$, or equivalently $n\ge d+\left\lceil \frac{d}{q}\right\rceil =g_q(2,d)$. As this lower bound applies to the nonlinear, systematic, and linear classes alike, $g_q(2,d)\le N_q(q^2,d)\le S_q(2,d)\le L_q(2,d)$. Conversely, for every $d$ there is a linear $[\,g_q(2,d),\,2,\,d\,]_q$ code \cite{SS65}, so $L_q(2,d)=g_q(2,d)$ and the three lengths coincide.  The length-maximal codes of Corollary~\ref{cor:2d_existence} realise the extremal cases $d=(s+1)q$.
	\end{proof}
	\begin{remark}
		The binary case of Corollary~\ref{cor:2d_LM_GB} was established in \cite{Guerrini2015Optimal}.
	\end{remark}

	\subsubsection{The Resolvable-Design Correspondence}
	\label{sec:resolvable_correspondence}

	The equality case of Lemma~\ref{lem:2d_bound} places us in classical territory. A length-maximal code of dimension two is an equidistant code with $q^2$ codewords meeting the Plotkin bound \eqref{eq:plotkin_length} with equality, and Semakov and Zinov'ev \cite{SemakovZinoviev1968} established that such codes are equivalent to resolvable balanced incomplete block designs. Proposition~\ref{prop:2d_design_equiv} below is the specialisation of their correspondence to the present setting, organised by the Singleton defect and stated for multidesigns (since repeated parallel classes arise as soon as $s \ge 1$). We include the short proof for completeness, as we shall make reference in the sequel.

	We allow repeated blocks throughout this subsection. A \emph{resolvable} $2$-$(v,K,\lambda)$ \emph{multidesign} is a multiset of $K$-subsets 	(\emph{blocks}) of a $v$-set in which every pair of points lies in exactly $\lambda$ blocks, counted with multiplicity, together with a partition of the block occurrences into \emph{parallel classes}, each class partitioning the point set \cite{BethJungnickelLenz1999}. Repetitions are necessary here because distinct code coordinates may induce the same partition. The letter $K$ denotes the block size, not the code dimension $k$.

	\begin{proposition}\label{prop:2d_design_equiv}
		For $q \ge 2$ and $s \ge 0$, the length-maximal $(n, q^2, d)_q$ $A^s$MDS codes correspond, up to code equivalence, to the resolvable $2$-$(q^2,\, q,\, s+1)$ multidesigns having $n$ parallel classes.
	\end{proposition}

	\begin{proof}
		Let $\mc{C}$ be such a code. By Lemma~\ref{lem:2d_bound}, $\mc{C}$ is symbol-uniform and equidistant, every pair of codewords agreeing in exactly $s+1$ coordinates. 
		Take the $q^2$ codewords as points. Each coordinate $i$ sorts the points into $q$ classes according to the symbol carried there. By symbol-uniformity each class has size $q$, so the coordinate $i$ provides a parallel class of $q$ blocks. 
		Two codewords agree in coordinate $i$ precisely when their corresponding points belong to a common block occurrence of that class. Equidistance implies that any two points lie together in exactly $s+1$ block occurrences. Hence we have a resolvable $2$-$(q^2,q,s+1)$ multidesign with one parallel class per coordinate. Conversely, labelling each point of such a multidesign by the block it occupies in each class yields a symbol-uniform, equidistant $(n,q^2,d)_q$ code with the stated parameters.
	\end{proof}

	The length formula $n = (s+1)(q+1)$ is, from this viewpoint, simply the design identity $r = \lambda(v-1)/(K-1)$ for the number of parallel classes: $r = (s+1)(q^2-1)/(q-1) = (s+1)(q+1)$, whence also $n \equiv s+1 \pmod q$.

	Proposition~\ref{prop:2d_design_equiv} converts the existence question for dimension-$2$ length-maximal codes into a question about resolvable multidesigns.

	\begin{corollary}\label{cor:2d_existence}
		Let $q \ge 2$ and $s \ge 0$.
		\begin{enumerate}
			\item For $s = 0$, a length-maximal code exists if and only if an affine plane of order $q$ exists. Moreover, the code is equivalent to a linear code if and only if the corresponding affine plane is Desarguesian. 
			\item If $q$ is a prime power, a length-maximal $(n, q^2, d)_q$ $A^s$MDS code exists for every $s \ge 0$.
			\item If a non-Desarguesian affine plane of order $q$ exists, then a \emph{genuinely nonlinear} length-maximal $(n, q^2, d)_q$ $A^s$MDS code exists for every $s \ge 0$. The least such $q$ is $9$.
		\end{enumerate}
	\end{corollary}

	\begin{proof}
%		(1) A resolvable $2$-$(q^2, q, 1)$ design is precisely an affine plane of order $q$. For the second part, and the more general concept of Desarguesian net-codes, see \cite{AldersonNetCodes2008}.\\
(1) A resolvable $2$-$(q^2,q,1)$ design is precisely an affine plane of order $q$, so the first assertion follows from Proposition~\ref{prop:2d_design_equiv}. Under this correspondence, the associated code is an index-$1$ net code based on that affine plane. By \cite[Theorem~2.5]{AldersonNetCodes2008}, such a code is equivalent to a linear code if and only if its underlying net is Desarguesian. Since the net has degree $q+1$, this is equivalent to the affine plane being isomorphic to $\mathrm{AG}(2,q)$.\\
		(2) For prime-power $q$ an affine plane of order $q$ exists, giving a resolvable $2$-$(q^2, q, 1)$ design. Repeating each of its parallel classes $s+1$ times produces a resolvable $2$-$(q^2, q, s+1)$ multidesign, equivalent by Proposition~\ref{prop:2d_design_equiv} to the required code. \\
		(3) Let $\pi$ be a non-Desarguesian affine plane of order $q$, and apply the construction of~(2), repeating each parallel class $s+1$ times. By Proposition~\ref{prop:2d_design_equiv}, the resulting multidesign corresponds to a length-maximal code $\mc C$. In the terminology of \cite{AldersonNetCodes2008}, $\mc C$ is a net code of uniform index $s+1$ based on $\pi$; repetition does not change its underlying net \cite[Lemma~2.1]{AldersonNetCodes2008}.
		
		If $\mc C$ were equivalent to a linear code, then it would be a Desarguesian net code by \cite[Theorem~2.5]{AldersonNetCodes2008}. Since equivalent net codes are based on isomorphic nets \cite[Theorem~2.4]{AldersonNetCodes2008}, this would imply
		\[
		\pi\cong\mathrm{AG}(2,q),
		\]
		contrary to the choice of $\pi$. Hence $\mc C$ is not equivalent to a linear code.
		
		
%		(3) Apply the repetition construction of~(2) to a non-Desarguesian affine plane $\pi$ of order $q$; by Proposition~\ref{prop:2d_design_equiv} the resulting multidesign is a length-maximal code $\mc{C}$. Code equivalence permutes coordinates and relabels symbols within a coordinate, so it preserves the multiset of partitions of $\mc{C}$ induced by the coordinates; this multiset is therefore an invariant, and for $\mc{C}$ it consists of the $q+1$ parallel classes of $\pi$, each with multiplicity $s+1$.
%
%		Suppose $\mc{C}$ were equivalent to a linear $[n,2,d]_q$ code $\mc{C}'$. Each coordinate of $\mc{C}'$ is a nonzero functional on $\F_q^2$, and two coordinates induce the same partition exactly when their functionals are scalar multiples, that is, when they determine the same point of $\mathrm{PG}(1,q)$. Let $m_P$ denote the multiplicity of the point $P$. A nonzero codeword $c_x$ has weight $n - m_{P(x)}$, where $P(x)$ is the point whose functionals annihilate $x$. By Lemma~\ref{lem:2d_bound} the code is equidistant of weight $d = n-(s+1)$, so $m_P = s+1$ for every one of the $q+1$ points. Hence the partitions of $\mc{C}'$ are the $q+1$ parallel classes of $\mathrm{AG}(2,q)$, each with multiplicity $s+1$, and the invariant forces $\pi \cong \mathrm{AG}(2,q)$, contradicting the choice of $\pi$. 
		
		
		The least order of a non-Desarguesian affine plane is $9$. The first example was constructed by Veblen and Maclagan-Wedderburn \cite{VeblenMaclaganWedderburn1907}; minimality follows from the classification through order 8, completed by Hall, Swift, and Walker \cite{HallSwiftWalker1956}.
	\end{proof}

	By contrast, for $s \ge 1$ and $q$ \emph{not} a prime power the existence question is open.

	\begin{remark}\label{rem:2d_open}
		The smallest undecided case is $q = 6$, $s = 1$. A length-maximal $(14, 36, 12)_6$ code is equivalent to a resolvable $2$-$(36, 6, 2)$ multidesign, whose existence is not settled by the non-existence of an affine plane of order $6$. More generally, determining the non-prime-power orders $q$ for which such a resolvable multidesign exists would complete the existence theory of dimension-$2$ length-maximal codes. It is likewise natural to ask whether, for prime-power $q$ and $s \ge 1$, every length-maximal code must arise from an affine plane by the repetition construction of Corollary~\ref{cor:2d_existence}.
	\end{remark}

	The dimension-$2$ estimate is the first member of a general shortening argument.

	\subsection{The General Bound}
	\label{sec:general_bound}

	\begin{theorem}[Shortening--Plotkin bound]\label{thm:main_bound}
		If $\mathcal{C}$ is an $(n, q^k, d)_q$ $A^s$MDS code with $k \ge 2$,
		then
		\begin{equation}\label{eqn:maximal_arc_bound}
			n \le (s+1)(q+1) + k - 2.
		\end{equation}
	\end{theorem}
	
	\begin{proof}
		Apply Lemma~\ref{lem:shortening} successively $k-2$ times. This produces a code $\mathcal C'$ with parameters $(n-k+2,q^2,d')_q$, where $d'\ge d$. The  Plotkin bound \eqref{eq:plotkin_length} gives
		\[
		 n-k+2\ge \left\lceil d'\frac{q+1}{q}\right\rceil \ge d\frac{q+1}{q}.
		\]
		Substituting $d=n-k+1-s$ and rearranging yields $n\le (s+1)(q+1)+k-2$.
	\end{proof}
	
	\begin{definition}\label{def:length_maximal}
		A code attaining the bound in Theorem~\ref{thm:main_bound} is called \emph{length-maximal}.
	\end{definition}
	
	\begin{remark}\label{rem:projective_connection}
		For a linear $[n,k,d]_q$ $A^s$MDS code, the columns of a generator matrix form an $(n, k+s-1)$-arc in $\mathrm{PG}(k-1,q)$, and
		codes attaining $n = (s+1)(q+1)+k-2$ correspond precisely to \emph{maximal arcs}. This connection, and the terminology \emph{length-maximal}, were introduced in \cite{alderson2025projectivesystems}. For $k = 3$, the theorem of Ball, Blokhuis, and Mazzocca
		\cite{Ball1997MaximalArcs} implies that non-trivial linear length-maximal codes (those with $s < q-2$) do not exist for odd $q$; for even $q$
		they exist at least whenever $(s+2) \mid q$, via Denniston's constructions \cite{Denniston1969}.
	\end{remark}
	
	\begin{corollary} \label{cor:griesmer_d_small}
		Let $\mathcal{C}$ be an $(n,q^k,d)_q$ A$^s$MDS code. If $k=2$, or if $k\ge 3$ and $d\le q^2$, then $\mathcal{C}$ respects the Griesmer bound.
	\end{corollary}

	\begin{proof}
		From Theorem \ref{thm:main_bound}, $n\le (s+1)(q+1)+k-2$. Substituting $s+1=n-k-d+2$ yields 
		\begin{equation}\label{eqn:singleton_to_griesmer}
			n\ge d+\frac{d}{q} +k-2.
		\end{equation}
		As $n$ is an integer, \eqref{eqn:singleton_to_griesmer} gives $n\ge d+\lceil d/q\rceil+k-2$. If $k=2$, or if $k\ge 3$ and $d\le q^2$, then $\lceil d/q^i\rceil = 1$ for $2\le i\le k-1$, so $g_q(k,d)= d+\lceil d/q\rceil +k-2$, and hence $n\ge g_q(k,d)$, the Griesmer bound \eqref{eqn:GriesmerBound}.
	\end{proof}
	
	\begin{remark}
		Corollary~\ref{cor:griesmer_d_small} subsumes Theorem~4 of \cite{Guerrini2015Optimal}, which shows that systematic codes with $d\le 2q$ respect the Griesmer bound.
	\end{remark}
	
	\begin{corollary} \label{cor:griesmer_s_small}
		Let $\mathcal{C}$ be an $(n,q^k,d)_q$ A$^s$MDS code, $k\ge 2$. If $s\le  q-1$, then $\mathcal{C}$ respects the Griesmer bound.
	\end{corollary}
	\begin{proof}
		By definition, $d=n-k+1-s$, so Theorem  \ref{thm:main_bound} provides   \begin{align*}
			d=n-k+1-s & \le (s+1)(q+1)+k-2-k+1-s\\ &= (s+1)(q+1) - (s+1) = q(s+1)\le q^2,
		\end{align*}  
		and the result follows from Corollary~\ref{cor:griesmer_d_small}.
	\end{proof}

	The following lemma provides a key consequence of a code being sufficiently long, that shortening preserves defect. It will be used repeatedly in the sequel.
	
	\begin{lemma}[Defect Preservation]\label{lem:defect_preserved}
		Let $\mathcal{C}$ be an $(n, q^k, d)_q$ $A^s$MDS code with $k \ge 3$, and let $\mathcal{C}'$ be an $(n-1, q^{k-1}, d')_q$ $A^{s'}$MDS code contained in a shortened code $\mathcal{C}_i^{(a)}$ (in particular, $\mathcal{C}'$ may be $\mathcal{C}_i^{(a)}$). If \[
		n \ge s(q+1) + k - 1,
		\]
		then $s' = s$ and $d' = d$.
	\end{lemma}

	\begin{proof}
		Since $\mathcal{C}'$ inherits minimum distance $d' \ge d$ from $\mathcal{C}$, we have $s' = n - k + 1 - d' \le s$. Suppose for contradiction that $s' \le s-1$. Applying Theorem~\ref{thm:main_bound} to $\mathcal{C}'$ gives
		\[
		n - 1 \le (s'+1)(q+1) + (k-1) - 2 \le s(q+1) + k - 3,
		\]
		so $n \le s(q+1)+k-2$, contradicting the hypothesis. Hence $s'=s$, and $d' = (n-1)-s-(k-1)+1 = n-s-k+1 = d$.
	\end{proof}
	
	The shortening--Plotkin bound has an equivalent minimum-distance formulation.
	
	\subsection{A Singleton-Type Reformulation}
	\label{sec:new_singleton}
	
	Substituting $s = n-k+1-d$ into \eqref{eqn:maximal_arc_bound} and isolating $d$ yields the following equivalent form. It may also obtained directly by shortening to $q^2$ codewords and applying the Plotkin bound.
	
	\begin{corollary}\label{cor:new_singleton}
		If $\mathcal{C}$ is an $(n, q^k, d)_q$ $A^s$MDS code with $k \ge 2$, then
		\[
		\left\lceil d \cdot \frac{q+1}{q} \right\rceil \le n - k + 2.
		\]
	\end{corollary}
	
	\begin{proof}
		From $n \le (s+1)(q+1) + k - 2$ and $s = n - k + 1 - d$, substitution gives $n \le (n - k + 2 - d)(q+1) + k - 2$, which provides $d \cdot \frac{q+1}{q} \le n - k + 2$.
	\end{proof}
	
	\begin{remark}
		Binary systematic codes satisfy $M = 2^k$, so Corollary~\ref{cor:new_singleton} gives $\lceil \frac{3}{2}d \rceil \le n - k + 2$, recovering \eqref{eq:GMS_bound} as a special case.
	\end{remark}
	
We now turn from length bounds to the internal structure of codes that achieve them.
	
	\subsection{Structural Consequences of Length-Maximality}
	\label{sec:structure}
	
	We now show that length-maximal codes are highly structured, in that they must be symbol-uniform, their pairwise distances take only a restricted set of values, and their defect satisfies a divisibility condition.
	
	\begin{theorem}[Symbol-uniformity under shortening]
		\label{thm:optimal_uniform}
		If $\mc{C}$ is a length-maximal $(n, q^k, d)_q$ $A^s$MDS code with $k \ge 2$, then $\mc{C}$ is symbol-uniform. Moreover, any code $\mc{C}'$ of dimension at least $2$
		obtained by successively shortening $\mc{C}$ is also symbol-uniform.
	\end{theorem}
	
	\begin{proof}
		We proceed by induction on $k$. The base case $k=2$ is the equality	characterisation in Lemma~\ref{lem:2d_bound}.
		
		Assume $k \ge 3$ and that $\mathcal{C}$ is length-maximal but not symbol-uniform. Some symbol $\alpha \in \A$ appears in at least $q^{k-1}+1$ codewords in some coordinate, say coordinate $1$. Consider the shortened code  $\mc{C}_1^{(\alpha)}$.	Since $|\mc{C}_1^{(\alpha)}| \ge q^{k-1}+1$, we may choose distinct	$u\ne  v \in \mc{C}_1^{(\alpha)}$.
		
		Since $n = (s+1)(q+1)+k-2 \ge s(q+1)+k-1$, Lemma~\ref{lem:defect_preserved} applies, so any $q^{k-1}$-subset $\mathcal{C}'\subseteq \mc{C}_1^{(\alpha)}$ has defect $s$, and is length-maximal. By the induction hypothesis, $\mc{C}'$  is symbol-uniform. Thus, by considering any two $q^{k-1}$-subsets of $\mc{C}_1^{(\alpha)}$ with symmetric difference $\{u, v\}$, it follows that $u$ and $v$ must carry the same symbol in each coordinate, contradicting $u \ne v$. Hence $\mathcal{C}$ is symbol-uniform.

		For the final claim, since $\mathcal{C}$ is symbol-uniform, every shortened code $\mathcal{C}_i^{(a)}$ has exactly $q^{k-1}$ codewords and, by Lemma~\ref{lem:defect_preserved}, defect $s$; its length is $(s+1)(q+1)+(k-1)-2$, so it is again length-maximal. Iterating, every code of dimension at least $2$ obtained from $\mathcal{C}$ by successive shortening is length-maximal, and is therefore symbol-uniform by the first part.
	\end{proof}

	\begin{corollary}[Orthogonal-array structure]\label{cor:orthogonal_array}
		Every length-maximal $(n,q^k,d)_q$ $A^s$MDS code with $k\ge 2$ is an orthogonal array
		\[
		\operatorname{OA}(q^k,n,q,k-1)
		\]
		of index $q$. More explicitly, if $1\le j\le k-1$, then any prescribed symbols in any $j$ prescribed coordinates occur together in exactly $q^{k-j}$ codewords.
	\end{corollary}

	\begin{proof}
		Fix $j$ coordinates and prescribed symbols. Shorten successively in the first $j-1$ of those coordinates. The resulting code has dimension $k-j+1\ge 2$ and is symbol-uniform by Theorem~\ref{thm:optimal_uniform}; hence the prescribed symbol in the last coordinate occurs $q^{k-j}$ times. This is the defining condition for an orthogonal array of strength $k-1$ and index $q$.
	\end{proof}
	
	\begin{theorem}[Distance Distribution of Length-Maximal Codes]
		\label{thm:optimal_weights}
		If $\mathcal{C}$ is a length-maximal $(n, q^k, d)_q$ $A^s$MDS code with $k \ge 2$, then for any two distinct codewords $u, v \in \mathcal{C}$,
		\[
		\dist(u,v) \in \{d\} \cup \{n, n-1, \ldots, n-k+3\}.
		\]
	\end{theorem}
	
	\begin{proof}
		We proceed by induction on $k$. For $k=2$, the result follows from Lemma~\ref{lem:2d_bound}.\\
		Assume $k \ge 3$. By Theorem~\ref{thm:optimal_uniform}, $\mathcal{C}$ is symbol-uniform. Let $u\ne v \in \mathcal{C}$. If $u$ and
		$v$ agree in no coordinate then $\dist(u,v) = n$. Otherwise, let $i$ be a coordinate in which $u_i = v_i$. The shortened code $
		\mathcal{C}_i^{(u_i)}$ contains both $u'$ and $v'$ (the shortening of $u$ and $v$ respectively). By 	Lemma~\ref{lem:defect_preserved}, $\mathcal{C}_i^{(u_i)}$ is a length-maximal $(n-1, q^{k-1}, d)_q$ $A^s$MDS code, and the induction hypothesis gives \[	\dist(u,v) = \dist(u',v') \in  \{d\} \cup \{n-1,\ \ldots,\ (n-1)-(k-1)+3\} = \{d\} \cup \{n-1, n-2, \ldots, n-k+3\},\]
		completing the induction.
	\end{proof}
	
	\begin{corollary}[Weight Distribution of Length-Maximal Codes]
		\label{cor:distance_distribution}
		If $\mathcal{C}$ is a length-maximal $(n, q^k, d)_q$ $A^s$MDS code with $k \ge 2$ and $\mathbf{0} \in \mathcal{C}$, then every nonzero codeword $\mathbf{c} \in \mathcal{C}$ has weight
		\[
		w(\mathbf{c}) \in \{d\} \cup \{n, n-1, \ldots, n-k+3\}.
		\]
	\end{corollary}

	\begin{proof}
		Apply Theorem~\ref{thm:optimal_weights} with $\mathbf{v} = \mathbf{0}$.
	\end{proof}

	\begin{lemma}\label{lem:divisibility}
		If $\mathcal{C}$ is a length-maximal $(n, q^k, d)_q$ $A^s$MDS code with $k \ge 3$, then $s \le q-1$ and $(s+2) \mid q(q+1)$.
	\end{lemma}
	
	\begin{proof}
		By Lemma~\ref{lem:defect_preserved}, we may successively shorten $\mathcal{C}$ to obtain a length-maximal $(n', q^3, d)_q$ $A^s$MDS code $\mathcal{C}'$ with $n' = (s+1)(q+1)+1$. We may assume (perhaps after relabelling) that $\mathbf{0} \in \mathcal{C}'$. By Corollary~\ref{cor:distance_distribution}, the nonzero codeword weights of $\mathcal{C}'$ lie in $\{d\} \cup \{n'\}$.
		
		Let $B$ denote the set of $q^2$ codewords of $\mathcal{C}'$ whose first entry is the symbol $1$. Since $\mathbf{0} \notin B$, each word in $B$ has weight $d$ or $n'$. A word of weight $d$ has
		$n' - d = s+2$ zeros, none in the first position. A word of weight $n'$ has no zeros. Thus every row of $B$ contributes either $0$ or $s+2$ zeros to columns	$2, \ldots, n'$.
		
		 The shortened code $\mathcal{C'}_1^{(1)}$ is a length-maximal $(n'-1, q^2, d)_q$ $A^s$MDS code  (Lemma~\ref{lem:defect_preserved}), so by Theorem~\ref{thm:optimal_uniform} is 	symbol-uniform. The total number of zero entries in $B$ is therefore $q \cdot (s+1)(q+1)$. Since each row contributes $0$ or $s+2$ zeros, $(s+2)$ divides $q(s+1)(q+1)$, and since	$\gcd(s+1, s+2) = 1$, we obtain $(s+2) \mid q(q+1)$.
		
		For the bound on $s$: the total zero count $q(s+1)(q+1)$ is at most  $q^2(s+2)$ (since each of the $q^2$ rows contributes at most $s+2$ zeros), so $q(s+1)(q+1) \le q^2(s+2)$, which simplifies to $s \le q-1$. 
			\end{proof}
		
	\begin{corollary}\label{cor:divisibility_gcd}
		If $\mathcal{C}$ is a length-maximal $(n, q^k, d)_q$ $A^s$MDS code 	with $k \ge 3$ and $\gcd(s+2, q) = 1$, then $(s+2) \mid (q+1)$. In particular, if $q$ is prime and $s < q-2$, then $(s+2) \mid (q+1)$.
	\end{corollary}
	
	\begin{proof}
		By Lemma~\ref{lem:divisibility}, $(s+2) \mid q(q+1)$; since $\gcd(s+2, q) = 1$, this forces $(s+2) \mid (q+1)$. If $q$ is prime and $s < q-2$, then $s+2 < q$ gives $\gcd(s+2, q) = 1$.
	\end{proof}
	
	\begin{remark}
		For linear length-maximal codes the stronger condition $s=q-1$ or $(s+2) \mid q$ is required \cite{Ball1997MaximalArcs, alderson2025projectivesystems}. 	Lemma~\ref{lem:divisibility} shows that nonlinear length-maximal codes must satisfy the weaker condition that $s\le q-1$ and $(s+2) \mid q(q+1)$, leaving open whether the stronger linear condition is also necessary in the nonlinear
		setting. Sufficient conditions under which $(s+2)\mid q$ is recovered are given in Lemma~\ref{lem:full_secant_divisibility} and Corollary~\ref{cor:homogeneous_divisibility} below.
	\end{remark}


	The divisibility of Lemma~\ref{lem:divisibility} is one facet of a sharper structural fact, that the zero-sets of the minimum-weight codewords form a $2$-design.

	\begin{theorem}[Minimum-weight zero-sets form a $2$-design]
		\label{thm:zero_set_design}
		Let $\mathcal{C}$ be a length-maximal $(n, q^3, d)_q$ $A^s$MDS code with $\mathbf{0} \in \mathcal{C}$, so that $n = (s+1)(q+1)+1$, $d = (s+1)q$, and every nonzero codeword has weight $d$ or $n$. For a weight-$d$ codeword $\mathbf{c}$ write $Z(\mathbf{c}) = \{\, i : c_i = 0 \,\}$ for its zero-set, of size $s+2$. Then the family
		\[
			\bigl\{\, Z(\mathbf{c}) : \mathbf{c} \in \mathcal{C},\ w(\mathbf{c}) = d \,\bigr\},
		\]
		regarded as a collection of blocks on the point set $\{1, \dots, n\}$, is a $2$-$(n,\, s+2,\, q-1)$ design with replication number $r = q^2 - 1$ and $b = n(q^2-1)/(s+2)$ blocks (counted with multiplicity).
	\end{theorem}

	\begin{proof}
		By Corollary~\ref{cor:distance_distribution} (with $k=3$) every nonzero codeword has weight $d$ or $n$; a weight-$n$ codeword has no zero entry, so every codeword possessing a zero coordinate is either $\mathbf{0}$ or a weight-$d$ word, whose zero-set has size $n-d = s+2$.

		Fix two coordinates $i \ne j$ and put $S = \{\, \mathbf{c} \in \mathcal{C} : c_i = 0 \,\}$. By symbol-uniformity (Theorem~\ref{thm:optimal_uniform}), $|S| = q^3/q = q^2$. Deleting coordinate $i$ from the words of $S$ produces the shortened code $\mathcal{C}_i^{(0)}$, a length-maximal $(n-1, q^2, d)_q$ $A^s$MDS code (Lemma~\ref{lem:defect_preserved}), which is itself symbol-uniform by Theorem~\ref{thm:optimal_uniform}. Hence in coordinate $j$ the symbol $0$ occurs exactly $|S|/q = q$ times among the words of $S$; that is, exactly $q$ codewords satisfy $c_i = c_j = 0$. One of these is $\mathbf{0}$, and the remaining $q-1$ each possess a zero coordinate and so have weight $d$. Therefore every pair of coordinates lies in exactly $q-1$ blocks, and the family is a $2$-design with $\lambda = q-1$.

		Each block has size $s+2$, so the replication number is $r = \lambda(n-1)/\bigl((s+2)-1\bigr) = (q-1)(s+1)(q+1)/(s+1) = q^2-1$; equivalently, each coordinate carries a zero in exactly $q^2-1$ of the $q^3-1$ nonzero codewords. Counting incidences yields $b(s+2) = n(q^2-1)$.
	\end{proof}

%	\begin{remark}\label{rem:design_geometry}
%		Theorem~\ref{thm:zero_set_design} locates the gap between the nonlinear condition $(s+2) \mid q(q+1)$ and the linear condition $(s+2) \mid q$. In the linear case the weight-$d$ codewords are the $q-1$ nonzero scalar multiples of each functional vanishing on a secant of the associated maximal arc. %; the design is then $q-1$ copies of the line-design of the arc, and 
%		That $(s+2) \mid q$ follows by counting the pencil of lines through a point external to the arc. This pencil count is  not recorded by the $2$-design of Theorem~\ref{thm:zero_set_design}, whose parameters are satisfiable with $(s+2) \nmid q$. For $q = 5$, $s = 1$ the zero-sets would form a $2$-$(13, 3, 4)$ design, realised for instance by four copies of a Steiner triple system $\mathrm{STS}(13)$. The $2$-design therefore neither forces nor forbids $(s+2) \mid q$. Whether the divisibility persists for nonlinear length-maximal codes is governed by structure finer than the design, and remains open.
%	\end{remark}

	A single block of \emph{full} multiplicity $q-1$ in the design of Theorem~\ref{thm:zero_set_design}, or equivalently, a single support shared by $q-1$ minimum-weight codewords, forces the same divisibility conditions of linear codes, provided $s \le q-2$.

	\begin{lemma}\label{lem:full_secant_divisibility}
		Let $\mathcal{C}$ be a length-maximal $(n, q^3, d)_q$ $A^s$MDS code with $s \le q-2$. If $\mathcal{C}$ contains $q-1$ weight-$d$ codewords sharing a common support, then $(s+2) \mid q$.
	\end{lemma}

	\begin{proof}
		We may assume $\mathbf{0} \in \mathcal{C}$ and, by Theorem~\ref{thm:optimal_uniform}, that $\mathcal{C}$ is symbol-uniform. Here $n = (s+1)(q+1)+1$ and $d = (s+1)q$, so $n-d = s+2$; by Theorem~\ref{thm:optimal_weights} any two distinct codewords agree in exactly $0$ or $s+2$ coordinates. Write $a(\mathbf{u},\mathbf{v})$ for the number of coordinates in which $\mathbf{u}$ and $\mathbf{v}$ agree.

		The $q-1$ given codewords share a zero-set of size $n-d = s+2$, which we take to be $\{1,\dots,s+2\}$. Let $B$ be the set of all codewords vanishing on $\{1,\dots,s+2\}$. Such a codeword has at least $s+2$ zeros, so it is either $\mathbf{0}$ or a weight-$d$ codeword with zero-set \emph{exactly} $\{1,\dots,s+2\}$. Two of the latter already agree on $\{1,\dots,s+2\}$, so---agreements being $0$ or $s+2$---they differ in every one of their $(s+1)q$ support coordinates; as those entries are nonzero, at most $q-1$ such codewords exist. The hypothesis attains this bound, so $|B| = q$.

		Restricted to the columns $\{s+3,\dots,n\}$, the $q$ rows of $B$ pairwise differ in every column; with $q$ rows and $q$ symbols, each of these $(s+1)q$ columns is a permutation of the alphabet.

		Choose a codeword $x$ nonzero in each of the first $s+2$ coordinates (such an $x$ exists; see below). In columns $1,\dots,s+2$ every row of $B$ is $0$ while $x$ is not, so $x$ agrees with no row there; in each of the remaining $(s+1)q$ columns exactly one row of $B$ matches $x$. Hence
		\[
		\sum_{r \in B} a(x,r) \;=\; (s+1)q .
		\]
		Each summand is $0$ or $s+2$, so $(s+2) \mid (s+1)q$; as $\gcd(s+1,s+2)=1$, this gives $(s+2) \mid q$.

		Finally, such an $x$ exists. Let $A$ be the set of codewords having a zero somewhere in $\{1,\dots,s+2\}$. Each coordinate is zero in exactly $q^2$ codewords (symbol-uniformity), of which $q$ lie in $B$, so
		\[
		|A| \;\le\; |B| + \sum_{i=1}^{s+2}\bigl(q^2 - q\bigr) \;=\; q + (s+2)(q^2-q) \;\le\; q^3 - q^2 + q \;<\; q^3,
		\]
		the penultimate step using $s+2 \le q$. Hence some codeword avoids $A$.
	\end{proof}

	\begin{remark}
		The hypothesis $s \le q-2$ is indeed required. The simplex code $[q^2+q+1,3,q^2]_q$ is length-maximal with $s = q-1$ and $(s+2) = q+1 \nmid q$.
	\end{remark}

	Since a linear code is closed under scalar multiplication, the $q-1$ nonzero scalar multiples of a minimum-weight codeword always supply such a shared support. The same holds under the weaker hypothesis of scalar-invariance alone.

	\begin{corollary}\label{cor:homogeneous_divisibility}
		Let $\mathcal{C}$ be a length-maximal $(n, q^3, d)_q$ $A^s$MDS code over $\mathbb{F}_q$ with $s \le q-2$ that is invariant under scalar multiplication, i.e.\ $\alpha\,\mathcal{C} = \mathcal{C}$ for every $\alpha \in \mathbb{F}_q^{\ast}$. Then $(s+2) \mid q$.
	\end{corollary}

	\begin{proof}
		Let $c$ be a weight-$d$ codeword. Its orbit $\{\alpha c : \alpha \in \mathbb{F}_q^{\ast}\}$ consists of $q-1$ distinct codewords, each of weight $d$ and with the same support as $c$. Apply Lemma~\ref{lem:full_secant_divisibility}.
	\end{proof}

	\begin{remark}
%		Corollary~\ref{cor:homogeneous_divisibility} extends the divisibility condition for linear codes ($(s+2)\mid q$) to the strictly larger class of scalar-invariant codes. 
		Contrapositively, by Lemma~\ref{lem:full_secant_divisibility}, the condition $(s+2)\mid q$ can fail to hold only if no minimum-weight support is shared by as many as $q-1$ codewords, equivalently, only if the design of Theorem~\ref{thm:zero_set_design} has no block of full multiplicity $q-1$.
	\end{remark}

	\subsection{Dimensions Two and Three: Inner Distance Distributions}
	\label{sec:dim23}

	For a nonlinear code, Hamming weight depends on the choice of a distinguished zero word. The invariant object is the inner distance distribution about a codeword $x\in\mathcal C$,
	\[
	A_i(x):=\bigl|\{c\in\mathcal C:\dist(x,c)=i\}\bigr|.
	\]
	Independent symbol permutations in the coordinates may be used to send any fixed $x$ to $\mathbf0$; this preserves all distances and length-maximality. Thus the pointed Hamming weight enumerator of the normalised code is exactly the inner distance distribution about $x$. The calculations below show, in particular, that this distribution is independent of $x$. 

	\begin{proposition}[Minimum-distance count]\label{prop:min_weight_count}
		Let $\mathcal{C}$ be a length-maximal $(n, q^k, d)_q$ $A^s$MDS code with $k \ge 2$, and set $t = n-d = s+k-1$. For every $x\in\mathcal C$, the number of codewords at distance $d$ from $x$ is
		\[
		A_d(x) = (q^2-1)\,\frac{\binom{n}{k-2}}{\binom{t}{k-2}}.
		\]
	\end{proposition}

	\begin{proof}
		Normalize the code by sending the chosen word $x$ to $\mathbf0$. By Theorem~\ref{thm:optimal_weights} the zero-count $z(\mathbf c) = n-w(\mathbf c)$ of a nonzero codeword lies in $\{0,1,\dots,k-3\}\cup\{t\}$, the value $t$ occurring exactly at distance $d$ from $x$. Corollary~\ref{cor:orthogonal_array} gives exactly $q^{k-j}$ codewords that are zero in any prescribed $j\le k-2$ coordinates. Therefore
		\[
		\sum_{\mathbf c \ne \mathbf 0} \binom{z(\mathbf c)}{j} = \binom{n}{j}\bigl(q^{k-j}-1\bigr).
		\]
		Take $j=k-2$. Only the case $z=t$ contributes on the left, giving $\binom{t}{k-2}A_d(x)=\binom{n}{k-2}(q^2-1)$.
	\end{proof}

	\begin{corollary}[Dimensions two and three]\label{cor:dim23_enumerator}
		Let $\mathcal{C}$ be a length-maximal $(n, q^k, d)_q$ $A^s$MDS code. For every baseword $x\in\mathcal C$, normalise $x$ to $\mathbf0$ and write $A_i=A_i(x)$.
		\begin{enumerate}
			\item If $k = 2$, then $\mathcal{C}$ is a one-weight code: with $n = (s+1)(q+1)$ and $d = (s+1)q$,
			\[
			A_0 = 1, \qquad A_d = q^2 - 1.
			\]
			\item If $k = 3$, then $\mathcal{C}$ is a two-weight code: with $n = (s+1)(q+1)+1$ and $d = (s+1)q$,
			\[
			A_d = \frac{n(q^2-1)}{s+2}, \qquad A_n = q^3 - 1 - A_d,
			\]
			and every other nonzero weight is absent.
		\end{enumerate}
	\end{corollary}

	\begin{proof}
		In each case Theorem~\ref{thm:optimal_weights} limits the nonzero weights, and Proposition~\ref{prop:min_weight_count} supplies $A_d$. For $k = 2$ the sole nonzero weight is $d$, so $A_d = q^2-1$. For $k = 3$ the nonzero weights are $d$ and $n$, and $\binom{n}{1}/\binom{t}{1} = n/(s+2)$ gives $A_d$, so   $A_n = q^3 - 1 - A_d$.
	\end{proof}

	\begin{remark}
		The dimension-three count $A_d = n(q^2-1)/(s+2)$ is exactly the number $b$ of blocks, counted with multiplicity, of the zero-set $2$-design of Theorem~\ref{thm:zero_set_design}. The two-weight enumerator and the $2$-design therefore carry the same information.
	\end{remark}

	\begin{remark}\label{rem:s0_mds}
		At $s=0$ a length-maximal code is a $(q+k-1,q^k,q)_q$ MDS code. Its inner distance distribution is determined by $n$, $k$, and $q$ alone \cite{Bush1952,Alderson2005Extending,Alderson2020}; in orthogonal-array language, an MDS code is an $\operatorname{OA}(q^k,n,q,k)$ of index one.
	\end{remark}

	For $k \ge 4$ the intermediate classes $A_{n-1}, A_{n-2}, \dots$ may also be determined. We compute these next.

	\subsection{Dimension Four: Distance Distribution and Existence}
	\label{sec:dim4}

	When $k=4$ we may determine the entire inner distance distribution and indicate when full-distance codewords may occur and when no length-maximal code exists. Fix any baseword, normalise it to $\mathbf0$, and write $A_w$ for the number of codewords at distance $w$ from it. Then $n=(s+1)(q+1)+2$, $d=(s+1)q$, and (Lemma~\ref{lem:divisibility}) $s\le q-1$. Put $t=n-d=s+3$. By Theorem~\ref{thm:optimal_weights}, the nonzero distances are $d$, $n-1$, and $n$, corresponding after normalisation to $t$, $1$, and $0$ zeros.

	\begin{theorem}[$k=4$ inner distance distribution]\label{thm:dim4_enumerator}
		With the notation above,
		\[
		A_d = \frac{n(n-1)(q^2-1)}{(s+3)(s+2)}, \qquad A_{n-1} = n(q^3-1) - (s+3)\,A_d,
		\]
		and the number of full-weight codewords is
		\[
		A_n = \frac{q(q-1)\,(t-q-1)\bigl((q+1)t-(4q+2)\bigr)}{t} \;=\; \frac{q(q-1)\,(s+2-q)\bigl((q+1)s-(q-1)\bigr)}{s+3}.
		\]
	\end{theorem}

	\begin{proof}
		We may assume $\mathbf 0 \in \mathcal C$, and by Theorem~\ref{thm:optimal_uniform} that $\mathcal C$ is symbol-uniform. With reference to Corollary \ref{cor:distance_distribution}, for nonzero $\mathbf{c}\in \mathcal C$ put $z(\mathbf c) = n - w(\mathbf c) \in \{0,1,t\}$; the three weight classes have sizes $A_n$, $A_{n-1}$, and $A_d$, with $A_n + A_{n-1} + A_d = q^4-1$.

		By symbol-uniformity each coordinate carries the symbol $0$ in exactly $q^3$ codewords, hence in $q^3-1$ nonzero codewords. Summing over the $n$ coordinates,
		 
		\begin{equation}
			\sum_{\mathbf c \ne \mathbf 0} z(\mathbf c) \;=\; n(q^3-1) \;=\; A_{n-1} + t\,A_d. \label{eqn1}
		\end{equation}

		Proposition~\ref{prop:min_weight_count}, in the case $k=4$, gives the minimum-weight count
		\[
		A_d = (q^2-1)\,\frac{\binom n2}{\binom t2} = \frac{n(n-1)(q^2-1)}{(s+3)(s+2)}.
		\]
		Equation~\eqref{eqn1} then yields $A_{n-1} = n(q^3-1) - t\,A_d$, and the total count gives $A_n = q^4-1-A_d-A_{n-1}$. Substituting and simplifying (using $n = (q+1)(t-2)+2$) yields the closed form for $A_n$, whose two expressions agree via $t = s+3$.
	\end{proof}

	As $s$ ranges over the admissible values $1 \le s \le q-1$, the sign of $A_n$ changes exactly once, yielding the following.

	\begin{corollary}\label{cor:dim4_trichotomy}
		No length-maximal $(n,q^4,d)_q$ $A^s$MDS code exists with $1 \le s \le q-3$. Consequently, if $\mathcal C$ is a length-maximal $(n,q^4,d)_q$ $A^s$MDS code, then exactly one of the following holds.
		\begin{enumerate}
			\item $s=0$, $q>2$, and $36\mid q$. (For $q=2$ the value $s=0$ coincides with $s=q-2$ and falls under case~(2).)
			\item $s = q-2$, and $A_n = 0$: $\mathcal C$ has no full-weight codeword.
			\item $s = q-1$, $(q+2)\mid 36$,  $A_{n-1} = 0$ and $A_n = \dfrac{q^2(q-1)^2}{q+2} > 0$, so $\mathcal C$ is a two-weight code with full-weight codewords.
		\end{enumerate}
	\end{corollary}

	\begin{proof}
		By Lemma~\ref{lem:divisibility}, $s \le q-1$. For $s \ge 1$, the sign of $A_n$ (Theorem~\ref{thm:dim4_enumerator}) is that of $t-q-1 = s+2-q$, which is negative for $1 \le s \le q-3$; since $A_n \ge 0$, no code exists in that range, and only $s \in \{0,\, q-2,\, q-1\}$ remains. Case (2) is immediate from Theorem~\ref{thm:dim4_enumerator}. Case (1) is part~(2) of Theorem~\ref{thm:bruen_silverman} with $k=4$; note that for $q = 2$ the value $s=0$ falls under case (2). For $s = q-1$, substitution gives $A_n = q^2(q-1)^2/(q+2)$ and, from Theorem~\ref{thm:dim4_enumerator}, $A_{n-1} = 0$. Moreover, integrality of $A_n$ requires $(q+2)\mid q^2(q-1)^2$. Observing  $q^2(q-1)^2 \equiv 36 \pmod{q+2}$ completes the result.
	\end{proof}
	
	\begin{remark}
		For each prime power $q$, codes as in part~(2) of Corollary~\ref{cor:dim4_trichotomy} exist owing to ovoids. An ovoid in $\mathrm{PG}(3,q)$ is a set of $q^2+1$ points meeting every line in at most two points, an elliptic quadric is a classical example \cite{Hirschfeld1985}. Every plane meets an elliptic quadric in at most $q+1$ points. Hence the points of the quadric form the projective system of a linear length-maximal $(q^2+1,q^4,q^2-q)_q$ $A^{q-2}$MDS code.
	\end{remark}


	\begin{corollary}\label{cor:dim4_existence}
		For an integer $q > 2$ with $36 \nmid q$, a length-maximal $(n,q^4,d)_q$ $A^s$MDS code exists only if $s \in \{q-2,\, q-1\}$. In the case $s=q-1$ additionally $(q+2)\mid 36$.
	\end{corollary}

	\begin{proof}
		Immediate from Corollary~\ref{cor:dim4_trichotomy}.
	\end{proof}

	\begin{remark}\label{rem:dim4_additive}
		In the additive (in particular linear) setting, Alderson and Ball \cite{alderson2026setssubspacesrestrictedhyperplane} establish the geometric counterpart of Corollary~\ref{cor:dim4_trichotomy}: a length-maximal additive code of dimension $4$ over $\mathbb F_{q^h}$ with $q>2$ must satisfy $n-d (=s+3) = q^h+1$ and have no full-weight (zero-secant) codewords. At $h=1$, Corollary~\ref{cor:dim4_trichotomy} recovers most of this combinatorially, for all (not necessarily additive) codes: the defect is confined to $\{0,\,q-2,\,q-1\}$, and in the case $s = q-2$ (that is, $n-d = q+1$) there is no full-weight codeword. Only the exclusion of $s = q-1$ is particular to the additive setting (cf.\ Remark~\ref{rem:dim4_q1_candidates}).
	\end{remark}

	\begin{remark}\label{rem:dim4_q1_candidates}
		Case~(3) of Corollary~\ref{cor:dim4_trichotomy} is the most restrictive. A length-maximal code of dimension four with $s = q-1$ has parameters $(q^2+q+2,\, q^4,\, q^2)_q$ and the forced two-weight enumerator $A_{n-1} = 0$, $A_n = q^2(q-1)^2/(q+2)$, $A_d = q^4-1-A_n$, and it can exist only if $(q+2)\mid 36$. As $q \ge 2$ gives $q+2 \ge 4$, the admissible alphabet sizes are exactly $q \in \{2,4,7,10,16,34\}$ (Table~\ref{tab:dim4_candidates}). The case $q = 2$ is realised by the extended Hamming code $[8,4,4]_2$, whose single full-weight codeword is the all-ones word. For every $q > 2$, by contrast, Remark~\ref{rem:dim4_additive} forbids a full-weight codeword in any additive (hence any linear) length-maximal dimension-four code, and for $q \in \{10,34\}$ the alphabet cannot be a field. So each of the five candidates with $q > 2$ would be \emph{genuinely nonlinear}. They thus form a complete, finite list of targets for the existence problem central to this paper, the smallest being a $(22,256,16)_4$ code.
	\end{remark}

	\begin{table}[htbp]
		\centering
		\renewcommand{\arraystretch}{1.15}
		\caption{The full-weight dimension-four family ($s = q-1$): length-maximal $(q^2+q+2,\, q^4,\, q^2)_q$ codes, which can exist only when $(q+2)\mid 36$. For $q > 2$ any such code is genuinely nonlinear (Remark~\ref{rem:dim4_additive}).}
		\label{tab:dim4_candidates}
		\begin{tabular}{@{}rclrrl@{}}
			\toprule
			$q$ & $q+2$ & parameters & $A_d$ & $A_n$ & existence \\
			\midrule
			$2$  & $4$  & $(8,16,4)_2$            & $14$      & $1$     & extended Hamming $[8,4,4]_2$ \\
			$4$  & $6$  & $(22,256,16)_4$         & $231$     & $24$    & open \\
			$7$  & $9$  & $(58,2401,49)_7$        & $2204$    & $196$   & open \\
			$10$ & $12$ & $(112,10^4,100)_{10}$   & $9324$    & $675$   & open; non-prime-power $q$ \\
			$16$ & $18$ & $(274,16^4,256)_{16}$   & $62335$   & $3200$  & open \\
			$34$ & $36$ & $(1192,34^4,1156)_{34}$ & $1301366$ & $34969$ & open; non-prime-power $q$ \\
			\bottomrule
		\end{tabular}
	\end{table}

	\subsection{Dimension Five: Non-Existence}
	\label{sec:dim5}

	The dimension $k=4$ enumerator propagates under shortening to a non-existence theorem in dimension $k=5$. Throughout, $\mathcal C$ is a length-maximal $(n,q^5,d)_q$ $A^s$MDS code, so $n = (s+1)(q+1)+3$, $d = (s+1)q$, and (Lemma~\ref{lem:divisibility}) $s \le q-1$. Set $t := n-d = s+4$; by Theorem~\ref{thm:optimal_weights} the nonzero weights are $d, n-2, n-1, n$, that is, $t, 2, 1, 0$ zeros with respective counts $A_d, A_{n-2}, A_{n-1}, A_n$.

	\begin{lemma}[$k=5$ inner distance distribution]\label{lem:dim5_enumerator}
		With the notation above,
		\[
		A_d = \frac{n(n-1)(n-2)(q^2-1)}{t(t-1)(t-2)}, \qquad
		A_{n-2} = \binom n2(q^3-1) - \binom t2\,A_d,
		\]
		\[
		A_{n-1} = n(q^4-1) - 2A_{n-2} - t\,A_d, \qquad
		A_n = q^5-1 - A_{n-1} - A_{n-2} - A_d.
		\]
	\end{lemma}

	\begin{proof}
		As in Theorem~\ref{thm:dim4_enumerator}, take $\mathbf 0 \in \mathcal C$, assume $\mathcal C$ symbol-uniform, and write $z(\mathbf c) \in \{0,1,2,t\}$ for the number of zero coordinates in $\mathbf c$. Proposition~\ref{prop:min_weight_count}, in the case $k=5$, gives the minimum-weight count
		\[
		A_d = (q^2-1)\,\frac{\binom n3}{\binom t3} = \frac{n(n-1)(n-2)(q^2-1)}{t(t-1)(t-2)}.
		\]
		For the remaining classes, shortening at $j-1$ all-zero coordinates yields, by Lemma~\ref{lem:defect_preserved}, a length-maximal code of dimension $6-j \ge 3$, necessarily symbol-uniform (Theorem~\ref{thm:optimal_uniform}), so any $j$ coordinates are all-zero in exactly $q^{5-j}$ codewords; therefore
		\[
		\sum_{\mathbf c \ne \mathbf 0}\binom{z(\mathbf c)}{j} = \binom nj\,(q^{5-j}-1), \qquad j=1,2.
		\]
		The $j=2$ identity gives $A_{n-2}$, $j=1$ gives $A_{n-1}$, and the total count $q^5-1$ then provides $A_n$.
	\end{proof}

	\begin{lemma}\label{lem:dim5_critical}
		At the two largest admissible defects,
		\[
		A_n\big|_{s=q-2} = -\frac{q^3(q-1)^2(q-3)}{2(q+2)}, \qquad
		A_n\big|_{s=q-1} = -\frac{q^2(q-1)^2(q^2-q+2)}{q+3}.
		\]
		Hence $A_n|_{s=q-1} < 0$ for every $q \ge 2$, while $A_n|_{s=q-2} < 0$ for $q \ge 4$ and $= 0$ for $q = 3$.
	\end{lemma}

	\begin{proof}
		Both are direct substitutions into Lemma~\ref{lem:dim5_enumerator}. At $s=q-2$ we have $n = q^2+2$, $t = q+2$, giving $A_{n-1}=0$, $A_{n-2} = \tfrac12(q^2+2)(q^2+1)(q-1)$, and $A_d = (q^2+2)(q^2+1)q(q-1)/(q+2)$. As such,  $A_n = (q^5-1) - A_d - A_{n-2}$ collapses to the stated form, whose numerator is $-q^3(q-1)^2(q-3)$. \\
		At $s=q-1$ we have $n = q^2+q+3$, $t = q+3$, giving $A_{n-2}=0$ and $A_d = (q^2+q+3)(q^2+q+2)(q^2+q+1)(q-1)/\bigl((q+3)(q+2)\bigr)$. So  $A_n = (q^5-1) - A_{n-1} - A_d$ simplifies to the stated form. The observation $q^2-q+2 = (q-\tfrac12)^2 + \tfrac74>0$ gives the sign claims.
	\end{proof}
	
		\begin{theorem}\label{thm:dim5existence}
		If $\mathcal{C}$ is a length-maximal $(n,q^k,d)_q$ A$^s$MDS code with $k\ge 5$, then exactly one of the following holds.
		\begin{enumerate}
			\item $s=0$ and $q=2$; here $\mathcal{C}$ is the trivial $(k+1,2^k,2)_2$ MDS code, necessarily equivalent to linear.
			\item $s=0$ and $36\mid q$; here $\mathcal{C}$ is necessarily genuinely nonlinear (Remark~\ref{rem:bs_nonlinear}).
			\item $s=1$ and $q=3$.
		\end{enumerate}
	\end{theorem}
	\begin{proof}
		By Lemma~\ref{lem:defect_preserved}, successive shortening yields a length-maximal $4$-dimensional code of the same defect $s$, so Corollary~\ref{cor:dim4_trichotomy} gives $s\in \{0,\, q-2,\, q-1\}$. Shortening instead to dimension $5$ and applying Lemma~\ref{lem:dim5_critical}, we have $s\ne q-1$, and if $s=q-2\ge 1$ then $q = 3$. If $s=0$, case (1) of Corollary~\ref{cor:dim4_trichotomy} forces $q=2$ or $36\mid q$; in the former case $\mathcal{C}$ is a binary MDS code, hence equivalent to linear \cite[Theorem~4]{Alderson2020}, namely the even-weight code of the stated parameters.
	\end{proof}

	\begin{theorem}\label{thm:dim5}
		Let $q \ge 4$. Every length-maximal code of dimension $k \ge 5$ over an alphabet of size $q$ is MDS. Consequently, if $36 \nmid q$, then no such code exists in any dimension $k \ge 5$.
	\end{theorem}

	\begin{proof}
		Take first $k=5$. Shortening $\mathcal C$ at a zero coordinate gives, by Lemma~\ref{lem:defect_preserved}, a length-maximal $(n-1,q^4,d)_q$ code of the \emph{same} defect $s$. By Corollary~\ref{cor:dim4_trichotomy}, no length-maximal dimension-four code has defect $1 \le s \le q-3$; together with $s \le q-1$ (Lemma~\ref{lem:divisibility}), the defect of $\mathcal C$ lies in $\{0, q-2, q-1\}$. For $s = q-1$, Lemma~\ref{lem:dim5_critical} gives $A_n < 0$ (an absurdity). For $s = q-2$ with $q \ge 4$, likewise $A_n < 0$. Hence $s = 0$. For $k \ge 6$, the result follows inductively by shortening. This proves the first assertion.

		The second assertion follows from shortening and Corollary~\ref{cor:dim4_existence}.
	\end{proof}

	\begin{remark}\label{rem:dim5_ball}
		Theorem~\ref{thm:dim5} provides a fully nonlinear counterpart of the additive non-existence theorem of Alderson and Ball \cite{alderson2026setssubspacesrestrictedhyperplane}. The excluded boundary $q = 3$ is genuinely exceptional in that a length-maximal dimension-five code \emph{does} exist. At $q = 3$, $s = q-2 = 1$, the forced enumerator $A_n = A_{n-1} = 0$, $A_d = 132$, $A_{n-2} = 110$ is exactly the weight distribution of the two-weight $[11,5,6]_3$ \emph{dual ternary Golay code} (nonzero weights $6$ and $9$), the projective two-weight code underlying the Berlekamp--van Lint--Seidel strongly regular graph $\mathrm{SRG}(243,22,1,2)$ \cite{CalderbankKantor1986}; indeed, every code with these parameters is equivalent to it (Remark~\ref{rem:uniqueness}). This exception continues one dimension further in that the extended ternary Golay code $[12,6,6]_3$ (weights $\{6,9,12\}$) is also length-maximal. We shall see that dimension $6$ is the true limit for $q=3$ (Proposition~\ref{prop:q3_dimension}).
	\end{remark}

	\subsection{The Ternary Case and Completeness}
	\label{sec:ternary_completeness}

	For the exceptional ternary codes, we record the inner distance distribution when $k=6$ and show that $k=7$ is not possible.
	\begin{proposition}\label{prop:dim6_q3}
		A length-maximal code over $\mathbb{F}_3$ of dimension $6$ has defect $s=1$ and parameters $(12,3^6,6)_3$. About every baseword its inner distance distribution is
		\[
		A_0 = 1, \quad A_6 = 264, \quad A_9 = 440, \quad A_{12} = 24,
		\]
		all other $A_w = 0$. It is realised by the extended ternary Golay code.
	\end{proposition}

	\begin{proof}
		By shortening, Theorem \ref{thm:dim5existence} forces $s = 1$, giving $n = 12$, $d = 6$, and $t = n-d = 6$, with zero-classes $z(\mathbf c) \in \{0,1,2,3,6\}$. Proposition~\ref{prop:min_weight_count} in the case $k = 6$ gives $A_6 = (q^2-1)\binom{12}{4}/\binom{6}{4} = 8 \cdot 495/15 = 264$. The identities $\sum_{\mathbf c \ne \mathbf 0}\binom{z(\mathbf c)}{j} = \binom{12}{j}(3^{6-j}-1)$ for $j = 1,2,3$ then give $A_9 = 440$ and $A_{10} = A_{11} = 0$, and the total $3^6-1$ gives $A_{12} = 24$. These are the weights $\{6,9,12\}$ of the extended ternary Golay code.
	\end{proof}

	\begin{remark}\label{rem:uniqueness}
		The extremal examples above are unique. Delsarte and Goethals \cite{DelsarteGoethals1975} proved that any (possibly nonlinear) code with the parameters of a binary or ternary Golay code, or of an extension thereof, is equivalent to that Golay code; the same paper (Theorem~4.4) establishes that any $(11,3^5,6)_3$ code is equivalent to the \emph{expurgated} ternary Golay code, that is, to the dual ternary Golay code. Consequently every length-maximal $(12,3^6,6)_3$ code is equivalent to the extended ternary Golay code, and every length-maximal $(11,3^5,6)_3$ code is equivalent to the dual ternary Golay code.

		The length-maximal $(8,2^4,4)_2$ code in case~(3) of Corollary~\ref{cor:dim4_trichotomy} is likewise unique up to equivalence. Normalize at a codeword and puncture one coordinate. The result is a perfect $(7,2^4,3)_2$ code with distribution $1+7z^3+7z^4+z^7$. Its seven weight-$3$ supports form the unique Steiner triple system on seven points, the Fano plane; the weight-$4$ supports are their complements, and the remaining nonzero word has full support. Thus the punctured code is the binary Hamming code. To restore the deleted coordinate while raising every distance-$3$ pair to distance at least four, the new bits must alternate on the graph joining pairs at distance three. That graph is connected because the weight-$3$ words span the Hamming code. Hence the extension is unique up to complementing the new coordinate, and is the overall-parity extension. Finally, every binary MDS code is equivalent to a linear code \cite[Theorem~4]{Alderson2020}, so every length-maximal $(k+1,2^k,2)_2$ code is equivalent to the even-weight code.
	\end{remark}

	The same method employed for Theorem~\ref{thm:dim5} determines the $q=3$ case completely.

	\begin{proposition}\label{prop:q3_dimension}
		Over an alphabet of size $q=3$, no length-maximal code of dimension $k \ge 7$ exists. Therefore, the maximal dimension of a length-maximal code over an alphabet of size $3$ is $6$.
	\end{proposition}

	\begin{proof}
		By shortening, it suffices to consider $k=7$. Again by shortening, Theorem \ref{thm:dim5existence} forces $s = 1$, so $n = 13$, $d = 6$, and $t = n-d = 7$, with zero-classes $z(\mathbf c) \in \{0,1,2,3,4,7\}$. The identities $\sum_{\mathbf c \ne \mathbf 0}\binom{z(\mathbf c)}{j} = \binom{13}{j}(3^{7-j}-1)$ for $j = 1,\dots,5$, together with the total $3^7-1$, force
		\[
		A_6 = \tfrac{3432}{7}, \quad A_9 = 1430, \quad A_{10} = A_{11} = 0, \quad A_{12} = 312, \quad A_{13} = -\tfrac{324}{7}.
		\]
		As $A_6, A_{13} \notin \mathbb{Z}$ (and $A_{13} < 0$), no length-maximal code of dimension $7$ exists.
	\end{proof}

	Assembling Theorem~\ref{thm:dim5existence} with the uniqueness results of Remark~\ref{rem:uniqueness}, the existence question in dimensions five and above is settled whenever $36 \nmid q$: only linear examples occur.

	\begin{corollary}[Linearity in dimension at least five]\label{cor:dim5_linearity}
		Let $\mathcal{C}$ be a length-maximal $(n, q^k, d)_q$ $A^s$MDS code with $k \ge 5$ and $36 \nmid q$. Then $\mathcal{C}$ is equivalent to a linear code, namely the even-weight $[k+1, k, 2]_2$ code, the dual ternary Golay $[11,5,6]_3$ code, or the extended ternary Golay $[12,6,6]_3$ code. In particular, a genuinely nonlinear length-maximal code of dimension $k \ge 5$ must be an MDS code with $36 \mid q$.
	\end{corollary}

	\begin{proof}
		Since $36 \nmid q$, Theorem~\ref{thm:dim5existence} leaves two cases. If $s = 0$ and $q = 2$, then $\mathcal{C}$ is a binary MDS code, equivalent to linear by \cite[Theorem~4]{Alderson2020}, necessarily the even-weight code. If $s = 1$ and $q = 3$, then $k \le 6$ (Proposition~\ref{prop:q3_dimension}), so $\mathcal{C}$ has parameters $(11, 3^5, 6)_3$ or $(12, 3^6, 6)_3$, and Remark~\ref{rem:uniqueness} gives equivalence to the dual or extended ternary Golay code, respectively.
	\end{proof}

	Collecting the dimensional bounds and enumerators yields a complete description.

	\begin{proposition}\label{prop:completeness}
		The inner distance distribution about every baseword of a length-maximal $(n,q^k,d)_q$ code is uniquely determined by $q$, $k$, and $s$. Equivalently, after any baseword is normalised to $\mathbf0$, the resulting Hamming weight enumerator is parameter-determined. For $s=0$ this is the distance distribution of a $(q+k-1,q^k,q)_q$ MDS code \cite{Alderson2020}. For $s\ge1$ a length-maximal code exists only when $k\le4$, or $q=3$ and $k\le6$, and the distribution is the explicit one furnished by Corollary~\ref{cor:dim23_enumerator}, Theorem~\ref{thm:dim4_enumerator}, Lemma~\ref{lem:dim5_enumerator}, or Proposition~\ref{prop:dim6_q3}.
	\end{proposition}

	\begin{proof}
		Let $s \ge 1$. By Theorem~\ref{thm:optimal_weights}, the zero-count of a nonzero codeword lies in the set $\{0,1,\dots,k-3\} \cup \{t\}$, so there are $k-1$ weight classes. The associated linear system is triangular: the top moment $j = k-2$ isolates $A_d$ (Proposition~\ref{prop:min_weight_count}), each successive lower moment introduces exactly one new class, and the total $q^k-1$ fixes the full-weight count $A_n$; back-substitution therefore determines all class sizes uniquely. The dimension restrictions come from Theorem~\ref{thm:dim5} (for $q \ge 4$) and Proposition~\ref{prop:q3_dimension} (for $q = 3$). For $q = 2$, Corollary~\ref{cor:binary_refinement_general} of Section~\ref{sec:binary_refinements}, which is proved independently of the present section, rules out $s \ge 2$, as it gives $n \le 3s+k-1 < 3s+k+1 = (s+1)(q+1)+k-2$; and Lemma~\ref{lem:binary_amds_bounds} caps the dimension at $k \le 4$ when $s = 1$. For $s = 0$ the code is MDS, with weight distribution determined by $n$, $k$, and $q$ alone \cite{Alderson2020}.
	\end{proof}

	Table~\ref{tab:enumerators} collects the resulting distributions dimension by dimension.

	\begin{table}[htbp]
		\centering
		\renewcommand{\arraystretch}{1.6}
		\caption{The inner distance distribution about any baseword of a length-maximal $(n,q^k,d)_q$ code with defect $s\ge1$, where $n=(s+1)(q+1)+k-2$, $d=(s+1)q$, and $t=n-d=s+k-1$; after normalising the baseword to zero, this is the pointed Hamming weight enumerator. For $s=0$ see Remark~\ref{rem:s0_mds}. The minimum-distance count is $A_d=(q^2-1)\binom{n}{k-2}/\binom{t}{k-2}$ (Proposition~\ref{prop:min_weight_count}). In the $k=5$ row, negativity of $A_n$ for the admissible nonzero defects yields non-existence when $q\ge4$ (Theorem~\ref{thm:dim5}). The $k=6$ row concerns $q=3$, $s=1$ (Proposition~\ref{prop:dim6_q3}).}
		\label{tab:enumerators}
		\begin{tabular}{@{}clll@{}}
			\toprule
			$k$ & nonzero distances & $A_d$ & $A_n$ (distance-$n$ count) \\
			\midrule
			$2$ & $d$ & $q^2-1$ & $0$ \\
			$3$ & $d,\,n$ & $\dfrac{n(q^2-1)}{s+2}$ & $q^3-1-A_d$ \\
			$4$ & $d,\,n{-}1,\,n$ & $\dfrac{n(n-1)(q^2-1)}{(s+3)(s+2)}$ & $\dfrac{q^2(q-1)^2}{q+2}\ (s{=}q{-}1);\ \ 0\ (s{=}q{-}2)$ \\
			$5$ & $d,\,n{-}2,\,n{-}1,\,n$ & $\dfrac{n(n-1)(n-2)(q^2-1)}{(s+4)(s+3)(s+2)}$ & $0\ (q{=}3);\ \ {<}0\ (q{\ge}4)$ \\
			$6$ & $6,\,9,\,12$ & $(q^2-1)\dbinom{n}{4}\big/\dbinom{t}{4}$ & $24\ (q{=}3)$ \\
			$\ge 7$ & \multicolumn{3}{l}{no code with $s\ge1$ exists (Propositions~\ref{prop:q3_dimension}, \ref{prop:completeness})} \\
			\bottomrule
		\end{tabular}
	\end{table}

	In the next section we show that additional hypotheses such as a binary alphabet or a large defect each permit  tighter bounds.

	%==============================================================================
	% SECTION 4: REFINEMENTS AND APPLICATIONS
	%==============================================================================
	\section{Refinements and Applications}
	\label{sec:refinements}
	
	In this section we first obtain parity-sensitive binary length bounds, extending a result of Guerrini, Meneghetti, and Sala. We then give the full shortening--Plotkin hierarchy, whose members improve Theorem~\ref{thm:main_bound} as the defect grows.

	\subsection{Binary Refinements}
	\label{sec:binary_refinements}

	We first recall the Hamming (sphere-packing) bound, which will be useful in this section.
	
	\begin{theorem}[Hamming Bound \cite{Hamming1950, MacWilliamsSloane1977}]
		\label{thm:hamming}
		Let $\mathcal{C}$ be an $(n,M,d)_q$ code, where $t = \left\lfloor \frac{d-1}{2} \right\rfloor$ is the maximum number of errors the code can correct. Then
		\[
		M \leq \frac{q^n}{\sum_{i=0}^{t} \binom{n}{i} (q-1)^i}.
		\]
	\end{theorem}
	
	We now proceed with a preliminary observation  that sufficiently long binary codes are forced to be symbol-uniform.
	
	\begin{lemma}\label{lem:binary_symbol_uniform_bound}
		Let $\mathcal{C}$ be a binary $(n, 2^k, d)_2$ $A^s$MDS code with $k \ge 2$.
		If $n > 3s + k - 1$, then $\mathcal{C}$ is symbol-uniform.
	\end{lemma}
	\begin{proof}
		The base case $k = 2$: suppose $\mathcal{C}$ is not symbol-uniform. Some coordinate contains a symbol appearing in at least $3$ codewords; say
		$a_1 = b_1 = c_1 = 1$ for distinct $\mathbf{a}, \mathbf{b}, \mathbf{c} \in \mathcal{C}$. In each of the remaining $n-1$ coordinate positions at least two of these three codewords must agree, contributing at least one agreement to each position. Since any two codewords agree in at most $n - d = s+1$ positions, we get $n-1 \le 3(s+1)-3 = 3s$, contradicting $n > 3s+1$.
		
		Assume now that $k\ge3$ and that the result holds in dimension $k-1$. If $\mathcal C$ is not symbol-uniform, some coordinate has a fibre containing at least $2^{k-1}+1$ codewords. Any $2^{k-1}$-subset of this fibre, after deletion of the fixed coordinate, has defect $s$ by Lemma~\ref{lem:defect_preserved}; its length $n-1$ satisfies $n-1>3s+(k-1)-1$, so it is symbol-uniform by induction. Choose two such subsets whose symmetric difference is a pair $\{u,v\}$. Equality of their symbol counts in every remaining coordinate forces $u$ and $v$ to agree everywhere, a contradiction. Thus $\mathcal C$ is symbol-uniform.
	\end{proof}
	
	Before discussing further refinements, let us situate the binary case. For $s = 0$, the MDS case is completely settled, in that every binary MDS code is equivalent to linear \cite[Theorem~4]{Alderson2020}, and thus satisfies $n\le k+1$ (with equality achievable), and the weight spectra of binary MDS codes are determined explicitly. In particular, no genuinely nonlinear binary MDS codes exist, so binary MDS codes offer no room for nonlinear improvement.
	
	For binary codes with $s = 1$, linear length-maximal ($n=k+4$) codes exist if and only if $k \le 4$, and for $k \ge 5$ a linear binary AMDS code has maximum length $k+2$ (see e.g.\ Corollary~4.25 in \cite{alderson2025projectivesystems}). The nonlinear case mirrors that of the linear.

	\begin{lemma}\label{lem:binary_amds_bounds}
		A binary $(n, 2^k, d)_2$ AMDS code satisfies $n \le k+2$ for $k \ge 5$. In particular, length-maximal binary $A^1$MDS codes (of length $k+4$) exist
		if and only if $2 \le k \le 4$.
	\end{lemma}

	\begin{proof}
		Suppose first that $n = k+3$, so that $d = n-k = 3$. The Hamming bound (Theorem~\ref{thm:hamming}) gives
		\[ 2^k \le \frac{2^n}{\sum_{i=0}^{1} \binom{n}{i}} = \frac{2^{k+3}}{1 + (k+3)} =\frac{8 \cdot 2^k}{k+4}, \]
		which requires $k \le 4$. Hence no binary $(k+3, 2^k, 3)_2$ $A^1$MDS code exists for $k \ge 5$.

		This excludes \emph{all} lengths $n \ge k+3$: an $A^1$MDS code has $d = n-k \ge 3$, so puncturing a coordinate causes no collisions and yields a binary $(n-1, 2^k, d')_2$ code with $d' \ge d-1$; for $n-1 \ge k+2$ the value $d' = (n-1)-k+1$ would make it a binary MDS code of length exceeding $k+1$, which is impossible, so $d' = (n-1)-k$ and the code is again $A^1$MDS. Puncturing a code of length $n \ge k+3$ down to length $k+3$ therefore produces a $(k+3, 2^k, 3)_2$ $A^1$MDS code, a contradiction. Thus $n \le k+2$ for $k \ge 5$.

		Finally, for $2 \le k \le 4$ linear length-maximal examples of length $k+4$ exist (Corollary~4.25 of \cite{alderson2025projectivesystems}).
	\end{proof}
	The argument extends to $q = 3$, and more generally yields a bound in terms of the $q$-ary Hamming ball volume for every alphabet size $q$.
	
	\begin{lemma}\label{lem:qary_amds_hamming}
		Let $\mathcal{C}$ be an $(n, q^k, d)_q$ $A^1$MDS code with $k \ge 2$ and $n = k + 2q - 1$ (one step below length-maximal). Then
		\[
		V(k + 2q - 1,\, q-1) \le q^{2q-1},
		\]
		where $V(n, t) = \sum_{i=0}^{t} \binom{n}{i}(q-1)^i$ denotes the volume of the $q$-ary Hamming ball of radius $t$. Consequently, if $q = 2$, then $k \le 4$, and if $q = 3$, then $k \le 6$.
	\end{lemma}
	
	\begin{proof}
		Since $s = 1$ and $n = k + 2q - 1$, we have $d = n - k = 2q - 1$, so $t = \lfloor (d-1)/2 \rfloor = q - 1$. The Hamming bound (Theorem~\ref{thm:hamming}) gives
		\[
		q^k \le \frac{q^n}{V(n,\, q-1)} = \frac{q^{k+2q-1}}{V(k+2q-1,\, q-1)},
		\]
		which rearranges to $V(k+2q-1,\, q-1) \le q^{2q-1}$.
		
		For $q = 2$: $V(k+3, 1) = 1 + (k+3) \le 8$ gives $k \le 4$, recovering the binary case above.\\
		For $q = 3$: $	V(k+5, 2) = 1 + 2(k+5)^2 \le 3^5$ gives $k \le 6$, with equality at $k=6$.
	\end{proof}

	\begin{remark}\label{rem:hamming_dimension_caps}
		Lemma~\ref{lem:qary_amds_hamming} recovers the dimensional bounds on length-maximal $A^1$MDS codes for $q=2$ and $q=3$, as in Lemma~\ref{lem:binary_amds_bounds} and Proposition~\ref{prop:q3_dimension} respectively. Indeed, a length-maximal $A^1$MDS code has length $k+2q$ and minimum distance $2q$, and puncturing any coordinate yields a $(k+2q-1,\, q^k,\, d')_q$ code with $d' \ge 2q-1$, to which the sphere-packing computation in the proof of Lemma~\ref{lem:qary_amds_hamming} applies verbatim, since $\lfloor (d'-1)/2 \rfloor \ge q-1$.
	\end{remark}

	Let us now focus on the cases where $s\ge 2$.  First, we have the following refinement for 3-dimensional codes.
	
	\begin{lemma}\label{lem:binary_refinement}
		Let $j \ge 0$ be an integer. If $\mathcal{C}$ is a binary $(n, 2^3, d)_2$ $A^s$MDS code with $s \ge 3j+2$, then
		\[
		n \le 3s + 2 - 2j.
		\]
	\end{lemma}
	
	\begin{proof}
		Suppose to the contrary that $n \ge 3s+3-2j$. From $s = n - 3 + 1 - d$ we have $d = n-2-s$.

		If $d$ is odd, then $s \ge 3j+2$ provides
		\[
		2d+1 - n = n - 3 - 2s \ge (3s+3-2j) - 3 - 2s = s - 2j \ge j+2 > 0,
		\]
		so part~(2) of the binary Plotkin bound (Theorem~\ref{thm:plotkin_binary}) applies. Since
		\[
		\frac{d+1}{2d+1-n} = 1 + \frac{s+2}{n-3-2s} \le 1 + \frac{s+2}{s-2j} = 2 + \frac{2j+2}{s-2j} \le 2 + \frac{2j+2}{j+2} = 4 - \frac{2}{j+2} < 4,
		\]
		we obtain the contradiction $8 = |\mathcal{C}| \le 2\left\lfloor\frac{d+1}{2d+1-n}\right\rfloor \le 6$.

		If $d$ is even, then, as $d = n-2-s$ is odd when $n = 3s+3-2j$, necessarily $n \ge 3s+4-2j$, whence
		\[
		2d - n = n - 4 - 2s \ge s - 2j \ge j+2 > 0,
		\]
		and part~(1) of Theorem~\ref{thm:plotkin_binary} gives, exactly as above,
		\[
		8 = |\mathcal{C}| \le 2\left\lfloor \frac{d}{2d-n} \right\rfloor = 2\left\lfloor 1 + \frac{s+2}{n-4-2s} \right\rfloor \le 2\left\lfloor 2 + \frac{2j+2}{s-2j} \right\rfloor \le 6,
		\]
		again a contradiction. Hence $n \le 3s+2-2j$.
	\end{proof}

	By leveraging Lemma~\ref{lem:defect_preserved}, and shortening, Lemma~\ref{lem:binary_refinement} provides an extension to codes of dimension $k \ge 3$.
	
	\begin{corollary}\label{cor:binary_refinement_general}
		If $\mathcal{C}$ is a binary $(n, 2^k, d)_2$ $A^s$MDS code with $k \ge 3$ and $s \ge 2$, then
		\[
		n \le 3s + k - 1.
		\]
	\end{corollary}
	
	\begin{proof}
		Since $s \ge 2$, Lemma~\ref{lem:binary_refinement} applies with $j=0$, establishing the result for $k=3$. For $k>3$, suppose $n \ge 3s+k$.  By Lemma \ref{lem:shortening}, successive shortening may be performed to produce an $(n-(k-3), 2^3, d)_2$ $A^s$MDS code $\mathcal{C}'$, the defect $s$ being preserved since $n \ge 3s+k$ (Lemma~\ref{lem:defect_preserved}). Applying Lemma~\ref{lem:binary_refinement} to $\mathcal{C}'$ with $j=0$ gives $n-(k-3) \le 3s+2$, hence $n \le 3s+k-1$, a contradiction.
	\end{proof}
	
	\begin{example}\label{ex:NR}
		The Nordstrom--Robinson code $(16, 2^8, 6)_2$ 	(see \cite{NordstromRobinson1967, MacWilliamsSloane1977}) has $k=8$, $d=6$, and Singleton defect $s = 3$, and as such, this famous code once more plays an extremal role, meeting the bound of Corollary~\ref{cor:binary_refinement_general} with equality.
	\end{example}
	
	
	Using the bound in Corollary~\ref{cor:binary_refinement_general}, an argument entirely similar to that in the proof of Corollary~\ref{cor:griesmer_d_small} provides the following.
	\begin{corollary}\label{cor:GBbinary_s2}
	If $\mathcal{C}$ is a binary $(n, 2^k, d)_2$ $A^s$MDS code with $k \ge 3$, $s \ge 2$, and $d\le 8$, then $\mc{C}$ respects the Griesmer bound.
	\end{corollary}
	
	By Corollary~\ref{cor:griesmer_s_small}, binary codes with $s\le1$ respect the Griesmer bound. Pairing this with Corollary~\ref{cor:GBbinary_s2} gives the following.
	\begin{corollary}\label{cor:binary_d8_griesmer}
	 If $\mc{C}$ is an $(n, 2^k, d)_2$ $A^s$MDS code with $k\ge 3$, and  $d\le 8$, then $\mc{C}$ respects the Griesmer bound.
	\end{corollary}

	For large defect, we obtain the following stratified bound.	
	\begin{lemma}\label{lem:binary_refinement_k}
		Let $j \ge 0$ be an integer. If \(\mathcal{C}\) is a binary \((n, 2^k, d)_2\) $A^s$MDS code with $k\ge 3$ and \(s \ge 3j+2\), then
		\[
		n \le 3s - 2j+k-1.
		\]
	\end{lemma}
	\begin{proof}
		Since the asserted bound is decreasing in $j$, it suffices to prove it for $j$ maximal with $s \ge 3j+2$; we may therefore assume in addition that $s \le 3j+4$. By Corollary~\ref{cor:binary_refinement_general}, we need only consider $j\ge 1$. The base case $k=3$ is Lemma~\ref{lem:binary_refinement}. Assume $k \ge 4$ and suppose to the contrary that $n \ge 3s+k-2j$. By Lemma~\ref{lem:shortening}, there exists a shortened code $\mc{C}_i^{(a)}$, an $(n',2^{k'},d')_2$ A$^{s'}$MDS code with $n'=n-1$, $k'=k-1$, $d'\ge d$, and $s'\le s$. Let $\alpha = s - s' \ge 0$.

	Note that by assumption
	 \[n'=n-1\ge 3s+k-2j-1\ge 3(3j+2)+k-2j-1=k+5+7j\ge (k-1)+13=k'+13. \]
	 Theorem~\ref{thm:main_bound}, applied to the shortened code, gives $n'\le3s'+k'+1$, hence $s'\ge4$. If $s'=4$, Corollary~\ref{cor:binary_refinement_general} instead gives $n'\le3s'+k'-1=k'+11$, contradicting $n'\ge k'+13$. Therefore $s'\ge5$, and consequently $0\le\alpha=s-s'\le s-5$.

		If $\alpha = 0$, then $s' = s \ge 3j+2$, and the induction hypothesis 	gives $n-1 \le 3s+(k-1)-1-2j = 3s+k-2j-2$, so $n \le 3s+k-2j-1$, contradicting $n \ge 3s+k-2j$.

		If $\alpha \ge 1$, set $j' = j - \lceil\alpha/3\rceil$. Since $\alpha \le s-5 \le 3j+4-5 = 3j-1$, we have $\lceil\alpha/3\rceil \le j$, and $j' \ge 0$. Moreover,
		$s' = s-\alpha \ge 3j+2-\alpha \ge 3(j-\lceil\alpha/3\rceil)+2 = 3j'+2$, so the induction hypothesis applies, giving
		\[
		n-1 \le 3(s-\alpha)+(k-2)-2\!\left(j-
		\left\lceil\tfrac{\alpha}{3}\right\rceil\right).
		\]
		Hence $n \le 3s+k-2j-(3\alpha-2\lceil\alpha/3\rceil)-1 \le 3s+k-2j-2$, since $3\alpha-2\lceil\alpha/3\rceil \ge 1$ for $\alpha \ge 1$; this contradicts $n \ge 3s+k-2j$.
	\end{proof}
	
	The following corollary shows that the binary refinement implies the Griesmer bound when $d$ is small relative to $s$.
	
	\begin{corollary}[Griesmer Bound for Binary Codes with Large Defect]
		\label{cor:binary_griesmer_large_s}
		Let $\mathcal{C}$ be a binary $(n, 2^k, d)_2$ $A^s$MDS code with $k \ge 3$ and $s \ge 3j+2$ for some integer $j \ge 0$. If 	$d \le 2(j+1)$, then $\mathcal{C}$ satisfies the Griesmer bound.
	\end{corollary}
	
\begin{proof}
	By Lemma~\ref{lem:binary_refinement_k},
	\[
	n \le 3s + k - 1 - 2j.
	\]
	Substituting $s = n - k + 1 - d$ and rearranging gives
	\begin{equation}\label{eqn:binary_large_s_length}
	n \ge d + \left\lceil\frac{d}{2}\right\rceil + k - 1 + j.
	\end{equation}

	On the other hand, for $d \le 2(j+1)$,
	\begin{align}
		g_2(k,d) &= d + \left\lceil\frac{d}{2}\right\rceil + \sum_{i=2}^{k-1} \left\lceil\frac{d}{2^i}\right\rceil \notag\\
		&< d + \left\lceil\frac{d}{2}\right\rceil + (k-2) + \sum_{i=2}^\infty \frac{d}{2^i} \notag\\
		&= d + \left\lceil\frac{d}{2}\right\rceil + k - 2 + \frac{d}{2} \notag\\
		&\le d + \left\lceil\frac{d}{2}\right\rceil + k - 1 + j, \label{eqn:binary_large_s_griesmer}
	\end{align}
	since $\frac{d}{2} \le j+1$. Comparing \eqref{eqn:binary_large_s_length} and \eqref{eqn:binary_large_s_griesmer} completes the proof.
\end{proof}

	The binary refinements exploit parity of $d$, though we may still obtain similar results when the defect $s$ is large relative to $q>2$.
	
	\subsection{A Shortening--Plotkin Hierarchy}
	\label{sec:large_defect}
	
	 Theorem~\ref{thm:main_bound} employed shortening to $k=2$. Shortening instead to any intermediate dimension and applying the same Plotkin estimate gives a hierarchy of bounds.

	\begin{theorem}[Shortening--Plotkin hierarchy]\label{thm:plotkin_shortening_hierarchy}
		Let $\mathcal C$ be an $(n,q^k,d)_q$ $A^s$MDS code with $k\ge2$. For every integer $r$ with $2\le r\le k$,
		\begin{equation}\label{eq:plotkin_shortening_hierarchy}
		 n\le
		 \left\lfloor (s+r-1)\frac{q^r-1}{q^{r-1}-1}\right\rfloor+k-r.
		\end{equation}
	\end{theorem}

	\begin{proof}
		Shorten $k-r$ times to a code of length $N=n-k+r$, size $q^r$, and minimum distance $d'\ge d$. Its defect is $s'=N-r+1-d'=n-k+1-d'\le s$. The unconditional Plotkin inequality \eqref{eq:plotkin_length} gives
		\[
		 N\ge d'\frac{1-q^{-r}}{1-q^{-1}}
		   =(N-r+1-s')\frac{q^r-1}{q^{r-1}(q-1)}.
		\]
		Rearranging and using $s'\le s$ yields
		\[
		 N\le (s'+r-1)\frac{q^r-1}{q^{r-1}-1}
		 \le (s+r-1)\frac{q^r-1}{q^{r-1}-1}.
		\]
		Since $N$ is integral, adding $k-r$ proves \eqref{eq:plotkin_shortening_hierarchy}.
	\end{proof}

	The case $r=2$ is Theorem~\ref{thm:main_bound}; the case $r=3$ is the most convenient next member and  improves that bound when $s\ge q$. For fixed parameters one may take the minimum of \eqref{eq:plotkin_shortening_hierarchy} over all $2\le r\le k$; members with $r>3$ can give further improvements for larger defect.
	
	The bound below is the case $r=k=3$ of Theorem~\ref{thm:plotkin_shortening_hierarchy}. We  give a direct agreement-counting proof, parallel to that of Lemma~\ref{lem:2d_bound} since the equality analysis of each step is what yields the characterisation of the extremal codes provide in Proposition~\ref{prop:dim3_design_equiv}.

	\begin{lemma}[Dimension-three Plotkin bound]\label{lem:dim3_plotkin_sharp}
		Let $\mc{C}$ be an $(n, q^3, d)_q$ $A^s$MDS code. Then
		\[
		n \;\le\; (s+2)\,\frac{q^2+q+1}{q+1} \;=\; (s+2)q+ \frac{s+2}{q+1}.
		\]
	\end{lemma}

	\begin{proof}
		Here $k = 3$, and $n-d = s+2$. For each coordinate $i$ and symbol $a\in\A$, let $f_{i,a} = |\{\mathbf{c}\in\mc{C} : c_i = a\}|$, and let $A = \sum_{i=1}^{n}\sum_{a\in\A}\binom{f_{i,a}}{2}$ count, over all coordinates, the agreeing pairs of codewords. Each of the $\binom{q^3}{2}$ pairs of distinct codewords agrees in at most $n-d = s+2$ coordinates, so $A \le \binom{q^3}{2}(s+2)$. On the other hand, $\sum_a f_{i,a} = q^3$, so by convexity of $\binom{\cdot}{2}$ each coordinate contributes $\sum_a \binom{f_{i,a}}{2}\ge q\binom{q^2}{2}$, minimised when $f_{i,a}=q^2$ for all $a$. Hence
		\[
		n\cdot q\binom{q^2}{2} \;\le\; A \;\le\; \binom{q^3}{2}(s+2),
		\qquad\text{i.e.}\qquad
		n\cdot\frac{q^3(q^2-1)}{2} \;\le\; \frac{q^3(q^3-1)}{2}\,(s+2).
		\]
		Dividing yields $n \le (s+2)\dfrac{q^3-1}{q^2-1} = (s+2)\dfrac{q^2+q+1}{q+1}$.
	\end{proof}
	
		The agreement count in Lemma~\ref{lem:dim3_plotkin_sharp} is tight precisely when the code carries a resolvable-design structure, exactly as in the dimension-$2$ case (Proposition~\ref{prop:2d_design_equiv}). Both equivalences are instances of the Semakov--Zinov'ev correspondence \cite{SemakovZinoviev1968}, here with $q^3$ codewords and blocks of size $q^2$.
	
	\begin{proposition}\label{prop:dim3_design_equiv}
		For $q \ge 2$ and $s \ge 0$, an $(n, q^3, d)_q$ $A^s$MDS code $\mathcal{C}$ attains the bound of Lemma~\ref{lem:dim3_plotkin_sharp} if and only if $\mathcal{C}$ is symbol-uniform and equidistant. Such a code corresponds, up to code equivalence, to a resolvable $2$-$(q^3,\,q^2,\,s+2)$ multidesign with $n$ parallel classes. In particular, attainment requires $(q+1) \mid (s+2)$, in which case, writing $s+2 = (q+1)m$, the parameters take the form $n = m(q^2+q+1)$, $d = mq^2$.
	\end{proposition}
	
	\begin{proof}
		Equality in the chain $n\cdot q\binom{q^2}{2}\le A\le\binom{q^3}{2}(s+2)$ of Lemma~\ref{lem:dim3_plotkin_sharp} forces equality in both steps. The left inequality is equality iff every coordinate is balanced, $f_{i,a}=q^2$ for all $a$ (symbol-uniformity); the right inequality is equality iff every pair of distinct codewords agrees in exactly $n-d=s+2$ coordinates (equidistance). Given both, take the $q^3$ codewords as points; by symbol-uniformity each coordinate $i$ partitions the points into $q$ blocks of size $q^2$ according to the symbol carried there, a parallel class, and two codewords share a block of this class iff they agree in coordinate $i$. Equidistance says any two points lie together in exactly $s+2$ block occurrences, so the $n$ classes form a resolvable $2$-$(q^3,q^2,s+2)$ multidesign \cite{BethJungnickelLenz1999}; the converse is immediate. The replication number is $r = \lambda(v-1)/(K-1) = (s+2)(q^3-1)/(q^2-1) = n$, and $\gcd(q+1,\,q^2+q+1)=1$ forces $(q+1)\mid(s+2)$ for $n$ to be an integer.
	\end{proof}
	
	The $r=3$ member of Theorem~\ref{thm:plotkin_shortening_hierarchy} is the following convenient closed form.

	\begin{corollary}\label{cor:dim3_plotkin_floor}
		If $\mc{C}$ is an $(n, q^k, d)_q$ $A^s$MDS code with $k \ge 3$, then
		\[
		n \;\le\; \left\lfloor (s+2)\,\frac{q^2+q+1}{q+1}\right\rfloor + k - 3
		\;=\; q(s+2) + \left\lfloor\frac{s+2}{q+1}\right\rfloor + k - 3.
		\]
	\end{corollary}

	\begin{proof}
		Take $r=3$ in Theorem~\ref{thm:plotkin_shortening_hierarchy} and use
		$(q^3-1)/(q^2-1)=(q^2+q+1)/(q+1)=q+1/(q+1)$.
	\end{proof}

		The bound in Corollary \ref{cor:dim3_plotkin_floor} improves on Theorem~\ref{thm:main_bound} when $s \ge q$. Indeed, a direct computation gives
\[
\left[(s+2)\frac{q^2+q+1}{q+1} + k - 3\right] - \bigl[(s+1)(q+1)+k-2\bigr]
\;=\; \frac{q\,(q-1-s)}{q+1},
\]
which is negative precisely when $s \ge q$. The comparison is made between the unfloored quantities; since the bound of Theorem~\ref{thm:main_bound} is already an integer, applying the floor in Corollary~\ref{cor:dim3_plotkin_floor} only sharpens the improvement.

For large defect we are able to provide a stratified bound.

    \begin{corollary}\label{cor:large_defect_bound_stratified}
    	If $\mc{C}$ is an $(n, q^k, d)_q$ $A^s$MDS code with $k \ge 3$, and $\beta(q+1)\le s+2 <(\beta+1)(q+1)$ for some integer $\beta\ge 1$, say $s+2=\beta(q+1)+r$ with $0\le r\le q$, then
    	\[
    	n\le (s+2-\beta)(q+1)+k-3+\beta-r.
    	\]
    	In particular, if $s\ge 2q$, then $n\le s(q+1)+k-1$.
   \end{corollary}
	
	\begin{proof}
		From Corollary \ref{cor:dim3_plotkin_floor}, we have
		\begin{align*}
			n& \le (s+2)q + \left\lfloor\frac{s+2}{q+1}\right\rfloor + k - 3\\
			& = (s+2)(q+1) -(s+2) + \left\lfloor\frac{s+2}{q+1}\right\rfloor + k - 3\\
			& = (s+2)(q+1) -(\beta(q+1)+r)+\beta + k - 3\\
			& = (s+2-\beta)(q+1) -r +\beta + k - 3.
		\end{align*}
		For the final claim, if $s\ge 2q$ then $s+2\ge 2(q+1)$, so $\beta\ge 2$; subtracting this bound from $s(q+1)+k-1$ leaves $q(\beta-2)+r\ge 0$, whence $n\le s(q+1)+k-1$.
	\end{proof}


	Likewise, we are able to provide stratified ranges within which the Griesmer bound holds.

    \begin{corollary}\label{cor:large_defect_griesmer}
    	Let $\mathcal{C}$ be an $(n, q^k, d)_q$ $A^s$MDS code with $k \ge 3$ and $s+2 \ge \beta(q+1)$, where $\beta< q+1$. If  $d \le \beta q^2$, then $\mathcal{C}$ respects the Griesmer bound.\\
    	In particular, if $s\ge 2q$ and $d\le 2q^2$, then $n\le s(q+1)+k-1$ and $\mathcal{C}$ respects the Griesmer bound.
    \end{corollary}
    
    \begin{proof}
    	Set $\beta' = \lfloor (s+2)/(q+1) \rfloor \ge \beta$. Corollary~\ref{cor:large_defect_bound_stratified} gives $n \le (s+2-\beta')(q+1)+k-3+\beta'$, and since each unit increase of $\beta'$ decreases this bound by $q$, also
    	\[
    	n \le (s+2-\beta)(q+1)+k-3+\beta.
    	\]
    	Substitute $s = n-k+1-d$ and isolate $n$ to obtain
    	\[
    	n \ge d + \left\lceil\frac{d}{q}\right\rceil+ \beta + k - 3.
    	\]
    	For $d \le \beta q^2$, $\lceil d/q^2\rceil \le \beta$, and each term $\lceil d/q^i\rceil =1$ for $i \ge 3$, so
    	\[
    	g_q(k,d) = \sum_{i=0}^{k-1}\left\lceil\frac{d}{q^i}\right\rceil
    	\le d + \left\lceil\frac{d}{q}\right\rceil + \beta + (k-3)
    	 \le n.
    	\]
    	The second part follows by substituting $\beta=2$.
    	\qedhere
    \end{proof}


	This completes the development of bounds and structural results. We now summarise the contributions and provide some open problems.
	%==============================================================================
	% SECTION 5: CONCLUSION
	%==============================================================================
	\section{Conclusion}
	\label{sec:conclusion}
	
	We established that $n \le (s+1)(q+1) + k - 2$ 	for every $(n,q^k,d)_q$ $A^s$MDS code with $k\ge2$. More generally, Theorem~\ref{thm:plotkin_shortening_hierarchy} supplies one bound for each shortening depth $2\le r\le k$. The main contribution of the paper is the rigidity in the cases of equality. A length-maximal code is symbol-uniform under every shortening of dimension at least two (Theorem~\ref{thm:optimal_uniform}), hence is an $\operatorname{OA}(q^k,n,q,k-1)$ of index $q$ (Corollary~\ref{cor:orthogonal_array}). Its pairwise distances lie in $\{d\}\cup\{n-k+3,\dots,n\}$, and for $k\ge3$ its defect satisfies $(s+2)\mid q(q+1)$. In dimension two, equality is equivalent to a resolvable $2$-$(q^2,q,s+1)$ multidesign and exists for every $s$ when $q$ is a prime power.

	We determine the inner distance distribution about every codeword. In dimension four they give a trichotomy: no length-maximal code exists for $1\le s\le q-3$; at $s=q-2$ no pair is at distance $n$; and at $s=q-1$ such pairs are forced and integrality requires $(q+2)\mid36$. For $q\ge4$, every length-maximal code of dimension at least five is MDS, so none exists when $36\nmid q$. The exceptional non-MDS cases are the dual and extended ternary Golay codes, of dimensions five and six. 
	
	The Singleton-type reformulation recovers and extends a bound of Guerrini, Meneghetti, and Sala for binary systematic codes. Nonlinear codes in specified parameter ranges are forced to respect the Griesmer bound, as summarised in Table~\ref{tab:griesmer_conditions}.
	

		\begin{table}[H] 
		\centering
		\caption{Conditions under which an $(n,q^k,d)_q$ $A^s$MDS code satisfies the
			Griesmer bound $n \ge g_q(k,d) := \sum_{i=0}^{k-1}\lceil d/q^i \rceil$.
			No linearity or systematicity assumption is made unless stated.}
		\label{tab:griesmer_conditions}
		\renewcommand{\arraystretch}{1.3}
		\begin{tabular}{@{}llr@{}}
			\toprule
			\textbf{Condition} & \textbf{Hypothesis} & \textbf{Result} \\
			\midrule
			\multicolumn{3}{@{}l}{\textit{General $q$, integer dimension
					($M = q^k$, $k \ge 2$)}} \\[2pt]
			$k = 2$
			& any $d$
			& Cor.~\ref{cor:2d_LM_GB} \\
			$k \ge 3$,\; $d \le q^2$
			& any $s$
			& Cor.~\ref{cor:griesmer_d_small} \\
			$s \le q-1$
			& any $d$\,${}^{(a)}$
			& Cor.~\ref{cor:griesmer_s_small} \\
			$s+2 \ge \beta(q+1)$,\; $d \le \beta q^2$,\; $1 \le \beta \le q$
			& $k \ge 3$
			& Cor.~\ref{cor:large_defect_griesmer} \\[6pt]
			$d \le 2q$
			& systematic${}^{(b)}$
			& \cite{Guerrini2015Optimal} \\
			$q^{k-1} \mid d$
			& any code, $M \ge q^k$
			& \cite{Guerrini2015Optimal} \\
			\midrule
			\multicolumn{3}{@{}l}{\textit{Binary ($q = 2$), integer dimension
					($M = 2^k$, $k \ge 3$)}} \\[2pt]
			$d \le 8$
			& any $s$
			& Cor.~\ref{cor:binary_d8_griesmer} \\
			$s \ge 3j + 2$,\; $d \le 2(j+1)$,\; $j \ge 0$
			& parametric family
			& Cor.~\ref{cor:binary_griesmer_large_s} \\
			\bottomrule
		\end{tabular}
		\smallskip
		
		\begin{minipage}{0.92\linewidth}\footnotesize
			$(a)$~Subsumed by row 2: $s \le q-1$ forces $d \le q^2$ via
			Thm.~\ref{thm:main_bound}. Listed separately for reference.\\
			$(b)$~Also subsumed by row 2 ($d \le 2q \le q^2$ for $q \ge 2$).
			The result in  \cite{Guerrini2015Optimal} applies to the stricter systematic class; row 2 applies to
			all codes.
		\end{minipage}
	\end{table}

		\medskip
	
	The principal open problem raised by this work is the following.

	\medskip
	
	\noindent\textbf{1. The odd $q$ problem.} 
	For linear codes, the theorem of Ball, Blokhuis, and Mazzocca \cite{Ball1997MaximalArcs} shows that non-trivial length-maximal $A^s$MDS codes do not exist when $q$ is odd and $s<q-2$. The known proof uses finite geometry. Beyond the single example of a $(21,9^3,18)_9$ code arising from the perp-system of De Clerck, Delanote, Hamilton, and Mathon \cite{DeClerck2002PerpSystems}, it is unknown whether genuinely nonlinear length-maximal codes can exist when linear ones do not. Are there further examples? Are there examples with $k>3$?
		\medskip
	
	Further natural questions include:
	\medskip

	\noindent\textbf{2. Unique extendability.}
	A classical result of Barlotti \cite{Barlotti1956}, recently extended to all dimensions \cite{alderson2025projectivesystems}, shows that a linear $A^s$MDS code of length $(s+1)(q+1)+k-3$ admits a unique extension to a length-maximal code whenever $(s+2) \mid q$ and $s < q-2$. The analogous question for nonlinear codes with $(k,s)\ne (2,0)$ is entirely open. For  $k\ge 4$, does an $(n, q^k, d)_q$ $A^s$MDS code with $n = (s+1)(q+1)+k - 3$ and $(s+2) \mid q(q+1)$ admit a unique extension to a length-maximal code? The symbol-uniformity and weight distribution results presented here may prove useful in approaching this problem.

	\medskip
	
	\noindent\textbf{3. Sharpening the divisibility condition.}
	Lemma~\ref{lem:divisibility} establishes $(s+2) \mid q(q+1)$, whereas linear length-maximal codes satisfy the stronger condition that $s = q-1$ or $(s+2) \mid q$ \cite{Ball1997MaximalArcs, alderson2025projectivesystems}. The boundary case $s=q-1$ must be excluded here: the simplex code $[q^2+q+1,3,q^2]_q$ is length-maximal with $s = q-1$ and $(s+2) = q+1 \nmid q$. Whether $(s+2) \mid q$ holds for all length-maximal codes with $s \le q-2$ is open. According to Theorem~\ref{thm:zero_set_design} it is not forced by the $2$-design structure of the dimension-$3$ core; by Lemma~\ref{lem:full_secant_divisibility} it can fail to hold only if no minimum-weight support is shared by as many as $q-1$ codewords. Either resolution would be of interest.  A proof of $(s+2) \mid q$ would align the nonlinear theory with the linear one, while a length-maximal code with $s \le q-2$ and $(s+2) \nmid q$ would furnish a genuinely nonlinear length-maximal code in a dimension $k \ge 3$ for which no linear one exists.

	No length-maximal MDS code of dimension three exists over an odd alphabet. Among the remaining pairs $(q,s)$ with $1 \le s \le q-2$ and $(s+2) \mid q(q+1)$ but $(s+2) \nmid q$, the smallest is $q=5$, $s=1$. The smallest undecided instance is therefore a $(13, 125, 10)_5$ code.
	\medskip

	\noindent\textbf{4. AMDS codes and the even-distance case.}
	Lemmas~\ref{lem:binary_amds_bounds} and~\ref{lem:qary_amds_hamming} use the Hamming bound to show that near-length-maximal $A^1$MDS codes of length $k+2q-1$ exist only for $k \le 4$ when $q=2$ and $k \le 6$ when $q=3$. The even-distance case $n = k+2q-2$ is inaccessible to the Hamming bound and remains open for large $k$. More generally, determining the maximum dimension of a near-length-maximal $A^s$MDS code for $s \ge 2$ and general $q$ is open.

	\medskip
	
	\noindent\textbf{5. Existence of length-maximal nonlinear codes.}
	In \cite{alderson2025projectivesystems}, the maximum dimension $k$ of a linear length-maximal $A^s$MDS code with $q>3$ was conjectured to be at most $5$. This conjecture has recently been settled in the affirmative  \cite{alderson2026setssubspacesrestrictedhyperplane}. For nonlinear codes the analogous question is no longer entirely open: Theorem~\ref{thm:dim5} shows that for every $q \ge 4$ with $36 \nmid q$ no length-maximal code of dimension $k \ge 5$ exists (and for every $q \ge 4$ such a code must be MDS), so the dimension is capped at four. The boundary $q = 3$ is exceptional: length-maximal codes exist there up to dimension six---the dual and extended ternary Golay codes---and no further (Proposition~\ref{prop:q3_dimension}). Moreover, every length-maximal code of dimension $k \ge 5$ with $36 \nmid q$ is equivalent to a linear one (Corollary~\ref{cor:dim5_linearity}), so genuinely nonlinear length-maximal codes can occur only in dimension $k \le 4$ or in the open MDS case $36 \mid q$. What remains genuinely open is attainment in dimension four. Here the full-weight case ($s=q-1$) is especially clean: a code exists only for $(q+2)\mid 36$, leaving the six alphabet sizes $q\in\{2,4,7,10,16,34\}$, of which only $q=2$ is realised (the extended Hamming code) and every $q>2$ would be genuinely nonlinear (Remark~\ref{rem:dim4_q1_candidates}, Table~\ref{tab:dim4_candidates}). The smallest such candidate is the two-weight $(22,256,16)_4$ code.

	\medskip
	
	This work provides a combinatorial foundation for the study of equality in the shortening--Plotkin bound and isolates concrete existence problems where genuinely nonlinear behaviour could first occur.
	
	\bibliographystyle{plain}
	\bibliography{references_v3}

\end{document}
